%ID VD_1_2_2_2_03
\textbf{Bước 1 — Xét hệ số trước khi chia.} Với $yu_x=x^2y$, trên $y\ne0$ chia cho $y$ được $u_x=x^2$. Trên $y=0$ phương trình gốc là $0=0$.

\textbf{Bước 2 — Tích phân theo $x$.} Giữ $y$ cố định:
\[u(x,y)=\int x^2\,dx+C(y)=x^3/3+C(y).\]

\textbf{Bước 3 — Kéo dài và kiểm tra.} Nếu miền chứa trục $y=0$ và yêu cầu nghiệm $C^1$, giá trị trên trục phải là giới hạn của công thức ở hai phía; không gán độc lập một hàm khác trên trục. Với $C$ đủ trơn, $yu_x=yx^2$ đúng cả ở $y=0$.

%ID VD_1_2_2_2_04
\textbf{Bước 1 — Đưa về ODE.} Với $u_x=2u$, coi $y$ là tham số và viết $u_x-2u=0$.

\textbf{Bước 2 — Dùng thừa số tích phân.} $M=e^{-2x}$ và
\[(e^{-2x}u)_x=e^{-2x}(u_x-2u)=0.\]
Tích phân cho $e^{-2x}u=C(y)$, tức $u=e^{2x}C(y)$. Cách này giữ cả nghiệm $u=0$, không cần chia cho $u$.

\textbf{Bước 3 — Kiểm tra.} $u_x=2e^{2x}C(y)=2u$. Chưa có dữ liệu đầu nên $C$ tùy ý đủ trơn.

%ID VD_1_2_2_2_05
\textbf{Bước 1 — Nhận dạng ODE tuyến tính.} Phương trình $u_x+2xu=4xy$ có $P(x)=2x$, nên $M=e^{\int2x\,dx}=e^{x^2}$.

\textbf{Bước 2 — Nhân và tích phân.}
\[(e^{x^2}u)_x=4xy e^{x^2}.\]
Giữ $y$ cố định, dùng $d(e^{x^2})=2xe^{x^2}dx$:
\[\int4xy e^{x^2}dx=2y e^{x^2},\quad e^{x^2}u=2y e^{x^2}+C(y).\]
Suy ra $u=2y+e^{-x^2}C(y)$.

\textbf{Bước 3 — Kiểm tra.} $u_x=-2xe^{-x^2}C(y)$, nên $u_x+2xu=4xy$.

%ID VD_1_2_2_2_06
\textbf{Bước 1 — Gom đạo hàm hỗn hợp.} Với $yu_{xy}+2u_x=x$, đặt $v=yu_y+2u$. Vì $y$ không đổi khi lấy đạo hàm theo $x$, $v_x=yu_{xy}+2u_x=x$.

\textbf{Bước 2 — Tích phân theo $x$.} $v=x^2/2+A(y)$, tức $yu_y+2u=x^2/2+A(y)$.

\textbf{Bước 3 — Giải theo $y$ trên $y\ne0$.} Chia cho $y$ rồi dùng thừa số $y^2$:
\[(y^2u)_y=yx^2/2+yA(y).\]
Tích phân cho $y^2u=x^2y^2/4+\int yA(y)dy+C(x)$. Đặt $B(y)=y^{-2}\int yA(y)dy$, ta được
\[u=x^2/4+B(y)+C(x)/y^2.\]

\textbf{Bước 4 — Kiểm tra và trục suy biến.} $u_x=x/2+C'(x)/y^2$, $u_{xy}=-2C'(x)/y^3$, nên $yu_{xy}+2u_x=x$. Muốn kéo dài qua $y=0$, tổ hợp $B(y)+C(x)/y^2$ phải đủ trơn; không cho phép chọn tùy ý các hạng kỳ dị. Phương trình gốc trên trục yêu cầu $2u_x(x,0)=x$.

%ID BT_1_2_3_01
\textbf{Bước 1: xác định điều cần phân loại.} Theo phần lý thuyết trong sách, ``bậc'' của câu này được hiểu là \emph{cấp phương trình}: cấp cao nhất của đạo hàm riêng xuất hiện. Số mũ của đạo hàm và hàm như $\log u,e^u$ không làm tăng cấp đạo hàm.

\textbf{Bước 2: xét từng phương trình.}
(a) $u_{xx}+u_{yy}=0$: hai đạo hàm đều cấp 2, nên cấp 2.
(b) $u_{xxx}+u_{xy}+a(x)u_y+\log u=f(x,y)$: $u_{xxx}$ cấp 3, $u_{xy}$ cấp 2, $u_y$ cấp 1; $\log u$ cấp 0. Cấp cao nhất bằng 3. Nếu nghiệm thực, $\log u$ cần $u>0$ nhưng không thay đổi cấp.
(c) $u_x+cu_y=d$: các đạo hàm đều cấp 1, nên cấp 1.
(d) $uu_{xx}+u_y^2+e^u=0$: $u_{xx}$ cấp 2; $(u_y)^2$ vẫn chỉ là bình phương của đạo hàm cấp 1. Do đó cấp 2.

\textbf{Kết luận:} các cấp lần lượt là $\boxed{2,3,1,2}$.

%ID BT_1_2_3_02
\textbf{Bước 1: dùng định nghĩa theo cấp cao nhất $m$.} Tuyến tính: ẩn và mọi đạo hàm xuất hiện tuyến tính, hệ số chỉ phụ thuộc biến độc lập. Nửa tuyến tính: đạo hàm cấp $m$ tuyến tính với hệ số chỉ phụ thuộc biến độc lập, nhưng các hạng cấp thấp có thể phi tuyến. Tựa tuyến tính: đạo hàm cấp $m$ tuyến tính, hệ số có thể phụ thuộc $u$ và các đạo hàm cấp thấp hơn. Hoàn toàn phi tuyến: biểu thức không tuyến tính theo đạo hàm cấp $m$. Chỉ hỏi thuần nhất khi phương trình đã tuyến tính.

\textbf{Bước 2: áp dụng từng ý.}
(a) $e^{\Delta u}=x$. Vì $\Delta u=u_{xx}+u_{yy}$ nằm trong hàm mũ, biểu thức đã cho hoàn toàn phi tuyến cấp 2. Tuy nhiên, nghiệm thực chỉ có trên $x>0$, và ở đó lấy log được phương trình tương đương $\Delta u=\log x$, tuyến tính không thuần nhất. Cần phân biệt phân loại \emph{biểu thức được in} và dạng tương đương sau biến đổi.

(b) $u_t+u_{txx}+uu_x=0$. Đạo hàm cao nhất $u_{txx}$ cấp 3, tuyến tính với hệ số 1. Tích $uu_x$ phi tuyến ở cấp thấp. Kết luận: nửa tuyến tính cấp 3.

(c) $u_t+(f(u))_x=0$. Quy tắc dây chuyền cho $u_t+f'(u)u_x=0$. Hạng cấp 1 tuyến tính theo $u_t,u_x$ nhưng hệ số $f'(u)$ phụ thuộc $u$, nên tựa tuyến tính cấp 1. Ngoại lệ: nếu $f(u)=au+b$, được $u_t+au_x=0$, tuyến tính thuần nhất.

(d) $u_{xx}+u_{yy}-2u=x^2$. Các hạng của $u$ tuyến tính, hệ số hằng, vế phải khác 0: tuyến tính không thuần nhất cấp 2.

(e) $u_{xy}=u$. Chuyển vế thành $u_{xy}-u=0$: tuyến tính thuần nhất cấp 2. Việc vế phải ban đầu là $u$ không làm phương trình không thuần nhất; phải gom tất cả hạng chứa ẩn về cùng vế.

(f) $uu_x+xu_y=0$. Tuyến tính theo hai đạo hàm cấp 1, nhưng hệ số của $u_x$ là $u$. Kết luận: tựa tuyến tính cấp 1.

(g) $u_x(1+u_y)=u_{xx}$. Viết $u_{xx}-u_x-u_xu_y=0$. Đạo hàm cấp 2 $u_{xx}$ tuyến tính, hệ số 1; phi tuyến $u_xu_y$ chỉ ở cấp 1. Kết luận: nửa tuyến tính cấp 2.

(h) $(\sin u_x)u_x+u_y=e^x$. Đạo hàm cấp cao nhất là cấp 1 và $u_x\sin u_x$ không tuyến tính theo $u_x$. Kết luận: hoàn toàn phi tuyến cấp 1.

\textbf{Bước 3: kiểm tra cách phân biệt.} Các ý (b),(g) có phi tuyến nhưng chỉ ở cấp thấp, nên nửa tuyến tính. Các ý (c),(f) có hệ số đạo hàm cấp cao nhất phụ thuộc ẩn, nên tựa tuyến tính. Các ý (a),(h) phi tuyến ngay ở đạo hàm cấp cao nhất. Với các phương trình phi tuyến không dùng phân biệt ``thuần nhất/không thuần nhất'' theo định nghĩa toán tử tuyến tính.

%ID BT_1_2_3_03
\textbf{Bước 1: tính đạo hàm của nghiệm thử.} Với $u=e^{rx+sy}$,
\[u_x=ru,\quad u_y=su,\quad u_{xx}=r^2u,\quad u_{xy}=rsu,\quad u_{yy}=s^2u.\]
Vì $u\ne0$, sau khi thế vào PDE có thể chia cho $u$ để có phương trình đại số theo $r,s$. Đề dùng nghiệm dạng đặc biệt; một họ nghiệm mũ không tự động là nghiệm tổng quát của mọi PDE.

(a) $2u_x+3u_y-2u=0$ cho $2r+3s-2=0$. Với mỗi $s\in\mathbb R$, chọn $r=1-3s/2$, được một nghiệm $e^{(1-3s/2)x+sy}$. Để có nghiệm tổng quát, dùng bất biến $\xi=3x-2y$, đặt $u=e^xF(\xi)$. Tính đạo hàm cho $2u_x+3u_y=2u$, nên toàn bộ nghiệm là $u=e^xF(3x-2y)$.

(b) Vì PDE có ba biến, dùng $u=e^{rx+sy+tz}$. Khi đó $u_{xyz}=rst\,u$, nên điều kiện là $rst=1$. Chọn $r,s\ne0$ tùy ý và $t=1/(rs)$; ví dụ $u=e^{x+y+z}$. Tuyến tính cho phép cộng hữu hạn các mode này. Không thể bỏ biến $z$ vì nghiệm không phụ thuộc $z$ có $u_{xyz}=0$.

(c) $4u_{xx}-4u_{xy}+4u_{yy}=0$ cho $r^2-rs+s^2=0$. Viết $(r-s/2)^2+3s^2/4=0$: với $r,s$ thực chỉ có $r=s=0$, tức mode hằng. Với số phức, $r=(1\pm i\sqrt3)s/2$; phần thực và phần ảo của mode phức cho nghiệm thực. Nghiệm mũ thực không mô tả hết nghiệm điều hòa của phương trình elliptic này.

(d) $u_{xx}+u_{yy}=u$ cho $r^2+s^2=1$. Tham số hóa $r=\cos\theta,s=\sin\theta$, được $u_\theta=e^{x\cos\theta+y\sin\theta}$. Các ví dụ $e^x,e^y,e^{(x+y)/\sqrt2}$ đều đúng; cộng các nghiệm do tính tuyến tính vẫn cho nghiệm. Thế các đạo hàm ở bước 1 kiểm tra trực tiếp từng công thức.

\textbf{Chi tiết tìm họ nghiệm ở (a).} Đặc trưng là $X'=2,Y'=3,U'=2U$. Tích phân từ $(x_0,y_0,u_0)$ được $X=x_0+2s,Y=y_0+3s,U=u_0e^{2s}$. Vì $3X-2Y=3x_0-2y_0$ và $e^{-X}U=e^{-x_0}u_0$, đặt giá trị bất biến thứ hai là $F(3X-2Y)$. Đổi về $x,y$ được $u=e^xF(3x-2y)$. Kiểm tra bằng các đạo hàm $u_x=e^x(F+3F')$, $u_y=-2e^xF'$: $2u_x+3u_y-2u=e^x(2F+6F'-6F'-2F)=0$.

%ID BT_1_2_3_04
\textbf{Bước 1: tìm hai họ đặc trưng.} Nếu $a\ne0$, phương trình $am^2-bm+c=0$ có hai nghiệm thực phân biệt $m_1,m_2$ vì $\Delta=b^2-4ac>0$. Ta có $m_1+m_2=b/a$, $m_1m_2=c/a$.

\textbf{Bước 2: đổi biến.} Đặt $\xi=y-m_1x,\eta=y-m_2x,u=v(\xi,\eta)$. Jacobian bằng $m_2-m_1\ne0$, nên phép đổi không suy biến. Quy tắc dây chuyền:
\begin{align*}
u_{xx}&=m_1^2v_{\xi\xi}+2m_1m_2v_{\xi\eta}+m_2^2v_{\eta\eta},\\
u_{xy}&=-m_1v_{\xi\xi}-(m_1+m_2)v_{\xi\eta}-m_2v_{\eta\eta},\\
u_{yy}&=v_{\xi\xi}+2v_{\xi\eta}+v_{\eta\eta}.
\end{align*}
Hai hệ số của $v_{\xi\xi},v_{\eta\eta}$ bằng 0. Hệ số hỗn hợp bằng $2am_1m_2-b(m_1+m_2)+2c=-\Delta/a\ne0$. Do đó PDE tương đương $v_{\xi\eta}=0$.

\textbf{Bước 3: tích phân.} Từ $\partial_\eta(v_\xi)=0$, $v_\xi=A(\xi)$. Tích phân theo $\xi$ cho $v=F(\xi)+G(\eta)$. Suy
\[\boxed{u=F(y-m_1x)+G(y-m_2x).}\]
Đây là dạng yêu cầu với $\alpha=-m_1,\beta=1,\gamma=-m_2,\delta=1$.

\textbf{Trường hợp $a=0$.} Khi đó $b\ne0$, PDE là $\partial_y(bu_x+cu_y)=0$. Đặt $\xi=x,\eta=y-(c/b)x$, tính được $bu_{xy}+cu_{yy}=b v_{\xi\eta}$. Vẫn có $u=F(x)+G(y-(c/b)x)$, hai dạng tuyến tính độc lập.

\textbf{Kiểm tra:} thế đạo hàm của $F(y-m_ix)$ cho hệ số $am_i^2-bm_i+c=0$; từng hạng giải PDE, và tổng cũng giải do tuyến tính.

%ID BT_1_2_3_05
(a) Xem $y$ là tham số: $u_x=-u$. Nhân $e^x$ cho $(e^xu)_x=0$, nên $u=e^{-x}C(y)$. Kiểm tra $u_x+u=0$.

(b) Xem $x$ là tham số. ODE theo $y$ là $u_{yy}-x^2u=0$, có phương trình đặc trưng $\lambda^2-x^2=0$. Với $x\ne0$, $u=A(x)e^{xy}+B(x)e^{-xy}$. Với $x=0$, ODE thành $u_{yy}=0$, nên vết là affine theo $y$. Một dạng thống nhất, không mất nghiệm ở $x=0$, là
\[u=A(x)\cosh(xy)+B(x)\frac{\sinh(xy)}x,\]
trong đó $\sinh(xy)/x=y$ tại $x=0$ theo giới hạn. Hệ số $A,B$ đủ trơn cho lớp nghiệm mong muốn.

(c) Đặt $v=u_x$, được $v_x-v=0$, nên $v=B(y)e^x$. Tích phân theo $x$ cho $u=A(y)+B(y)e^x$. Đạo hàm lại cho $u_{xx}=u_x$.

(d) Với nghiệm thực, $u_x^2+u_y^2=0$ là tổng hai số không âm; buộc $u_x=u_y=0$. Do đó $u$ hằng trên từng thành phần liên thông. Nếu vế phải là $k\ne0$: $k<0$ thì vô nghiệm thực; $k>0$ là eikonal $|\nabla u|=\sqrt k$. Một họ nghiệm đơn giản $u=ax+by+c$, $a^2+b^2=k$, nhưng không phải toàn bộ nghiệm.

Để dựng nghiệm cục bộ tổng quát từ dữ liệu $u(\gamma(s))=h(s)$, chọn $P(s)$ thỏa $|P(s)|^2=k$ và $P(s)\cdot\gamma'(s)=h'(s)$. Hệ đặc trưng của $F=p^2+q^2-k$ có $P'=0$, $X'=2P$, $U'=2k$, nên
\[X=\gamma(s)+2tP(s),\qquad U=h(s)+2kt.\]
Khử $(s,t)$ ở nơi Jacobian của $X$ khác 0, được nghiệm $u(X)=U$. Đây là dạng áp dụng phương pháp đặc trưng, phân biệt với chỉ đưa ra ví dụ affine.

\textbf{Chi tiết xử lý điểm $x=0$ ở (b).} Chọn hai hệ số từ các vết tại $y=0$: $A(x)=u(x,0)$, $B(x)=u_y(x,0)$. Khi $x\ne0$, giải hệ $C+D=A$, $x(C-D)=B$ cho $C=(A+B/x)/2,D=(A-B/x)/2$, rồi gom hai hàm mũ thành $A\cosh(xy)+B\sinh(xy)/x$. Khai triển $\sinh(xy)/x=y+x^2y^3/3!+\cdots$ cho giới hạn $y$ và vết $u(0,y)=A(0)+B(0)y$; không loại nghiệm tuyến tính này khi hai nghiệm mũ trùng nhau.

%ID VD_1_3_1_2_01
\textbf{Bước 1 — Xác định toán tử.} Đặt $L=a\partial_x+b\partial_t$. Do lấy đạo hàm và nhân hệ số là các phép tuyến tính, với $c_1,c_2$ hằng,
\[L(c_1u_1+c_2u_2)=c_1Lu_1+c_2Lu_2.\]

\textbf{Bước 2 — Dùng giả thiết.} Nếu $Lu_1=Lu_2=0$, vế phải bằng 0, nên $c_1u_1+c_2u_2$ cũng là nghiệm. Điều này đúng cho mọi tổng hữu hạn các nghiệm.

\textbf{Bước 3 — Phạm vi chồng chất.} Với PDE có nguồn $Lu=f\ne0$, tổng hai nghiệm cùng nguồn có nguồn $2f$. Khi đó dùng $u=u_p+u_h$ với $Lu_p=f$, $Lu_h=0$; hiệu hai nghiệm cùng nguồn là nghiệm thuần nhất.

%ID VD_1_3_1_2_02
\textbf{Bước 1 — Đặc trưng.} Với $au_x+bu_t=0$ và $b\ne0$, dùng $t$ làm tham số: $X'=a/b$, $U'=0$. Do đó $X=\xi+at/b$, $u=F(x-at/b)$.

\textbf{Bước 2 — Áp các dữ liệu.} Với $u(x,0)=xe^{-x^2}$, $F(\xi)=\xi e^{-\xi^2}$, nên
\[u=(x-at/b)e^{-(x-at/b)^2}.\]
Với $u(0,t)=t$ và $a\ne0$, dùng $x$ làm tham số: $t-bx/a$ bất biến. Dữ liệu cho $u=t-bx/a$. Với $a=1,b=-2$ và dữ liệu $1/(1+x^2)$, có $u=1/[1+(x+t/2)^2]$.

\textbf{Bước 3 — Kiểm tra hệ số và ngoại lệ.} Với $u=F(x-at/b)$, $au_x+bu_t=aF'-aF'=0$. Với $u=t-bx/a$, $au_x+bu_t=-b+b=0$. Các hệ số $a,b$ trong hai kết quả in trong sách đã đổi vai trò; các công thức trên sửa theo PDE. Nếu hệ số đạo hàm pháp tuyến bằng 0, đường dữ liệu là đặc trưng và dữ liệu phải hằng dọc nó; các dữ liệu không hằng ở đây không phù hợp. Nếu $a=b=0$, PDE không ràng buộc nghiệm.

%ID BT_1_3_1_3_01
\textbf{Ý (a) — Đạo hàm theo hướng.} Nếu $(a,b)\ne(0,0)$, đặt $\nu=(a,b)/\sqrt{a^2+b^2}$. Khi đó $D_\nu u=(au_x+bu_y)/\sqrt{a^2+b^2}=0$. Nghiệm không đổi khi di chuyển theo hướng $(a,b)$.

\textbf{Ý (b) — Đặc trưng.} Giải $X'=a,Y'=b$: $(X,Y)=(x_0,y_0)+s(a,b)$. Tính $(bX-aY)'=ba-ab=0$, nên các đường là $bx-ay=\text{hằng}$.

\textbf{Ý (c) — Nghiệm.} $U'=au_x+bu_y=0$ trên mỗi đường, nên $u=F(bx-ay)$. Kiểm tra $u_x=bF'$, $u_y=-aF'$, suy $au_x+bu_y=0$. Nếu $a=b=0$, PDE là $0=0$, mọi hàm $C^1$ là nghiệm. Trên toàn mặt phẳng, mỗi mức của bất biến là một đường liên thông; trên miền bị chia cắt cần xét từng phần đặc trưng liên thông.

%ID BT_1_3_1_3_02
\textbf{Ý (a).} $u_x+u_y-u=0$. Bước 1: $dy/dx=1$ cho $\rho=y-x$. Bước 2: $U'=U$ cho $U=e^xF(\rho)$. Bước 3: $u_x=e^x(F-F')$, $u_y=e^xF'$, nên $u_x+u_y=u$.

\textbf{Ý (b).} $2u_t+3u_x=0$. Chia cho 2, dùng $t$ làm tham số: $X'=3/2$, $U'=0$. Bất biến $\rho=x-3t/2$, do đó $u=F(x-3t/2)$. Kiểm tra $2u_t+3u_x=2(-3F'/2)+3F'=0$.

\textbf{Ý (c).} $au_t+bu_x=u$ với $a,b\ne0$. $X'=b/a$ và $U'=U/a$ cho bất biến $\rho=ax-bt$, $U=e^{t/a}F(\rho)$. Do đó $u=e^{t/a}F(ax-bt)$. Đạo hàm cho $au_t=e^{t/a}(F-abF')$, $bu_x=ab e^{t/a}F'$, tổng bằng $u$. Các hàm $F$ tùy ý $C^1$.

%ID BT_1_3_1_3_03
\textbf{Ý (a) — $au_x+bu_y=f$, $a,b,f$ hằng.} Nếu $a\ne0$, dùng $x$ làm tham số: $dy/dx=b/a$, nên $\rho=bx-ay$. Trên đường, $dU/dx=f/a$, tích phân cho
\[u=fx/a+F(bx-ay).\]
Kiểm tra $au_x+bu_y=a(f/a+bF')-abF'=f$. Nếu $a=0,b\ne0$, tích phân theo $y$ cho $u=fy/b+F(x)$. Nếu $a=b=0$, chỉ có nghiệm khi $f=0$, khi ấy mọi hàm đủ trơn đều thỏa.

\textbf{Ý (b) — $u_x+2u_y+2u=1$.} Bước 1: $dy/dx=2$ cho $\rho=y-2x$. Bước 2: $U'+2U=1$. Nhân $e^{2x}$:
\[(e^{2x}U)'=e^{2x},\quad e^{2x}U=e^{2x}/2+F(\rho).\]
Suy $u=1/2+e^{-2x}F(y-2x)$. Bước 3: $u_x=e^{-2x}(-2F-2F')$, $u_y=e^{-2x}F'$, nên $u_x+2u_y+2u=1$.

\textbf{Ý (c) — $3u_x-4u_y=x+e^x$.} $dy/dx=-4/3$ cho $\rho=4x+3y$. Dọc đường, $U'=(x+e^x)/3$; tích phân:
\[u=x^2/6+e^x/3+F(4x+3y).\]
Kiểm tra $3u_x=x+e^x+12F'$, $-4u_y=-12F'$, tổng bằng vế phải.

\textbf{Ý (d) — $u_x+3u_y=9y^2$.} Bất biến $\rho=y-3x$. Trên đường $y=3x+\rho$, có $U'=9(3x+\rho)^2$. Vì $d[(3x+\rho)^3]/dx=9(3x+\rho)^2$, được $u=y^3+F(y-3x)$. Kiểm tra $u_x=-3F'$, $u_y=3y^2+F'$, tổng bằng $9y^2$.

\textbf{Ý (e) — Dùng $h(x)=2x^2+3x$ trên $y=0$.} Thực hiện bước xác định hàm tùy ý riêng cho từng PDE:

Với (a), nếu $b\ne0$, đặc trưng từ $(\xi,0)$ có $x=\xi+ay/b$. Do đó $\xi=x-ay/b$, còn $dU/dy=f/b$, $U(0)=h(\xi)$. Suy
\[u=h(x-ay/b)+fy/b.\]
Trên $y=0$ bằng $h(x)$; $au_x+bu_y=ah'(\xi)+b[-ah'(\xi)/b+f/b]=f$.

Với (b), $h(x)=1/2+e^{-2x}F(-2x)$. Đặt $s=-2x$, được $F(s)=e^{-s}[h(-s/2)-1/2]$. Thay $s=y-2x$:
\[u=\tfrac12+e^{-y}[h(x-y/2)-\tfrac12].\]

Với (c), đặt $P(x)=x^2/6+e^x/3$. Dữ liệu cho $F(4x)=h(x)-P(x)$, nên $F(s)=h(s/4)-P(s/4)$. Suy
\[u=P(x)+h(x+3y/4)-P(x+3y/4).\]

Với (d), $F(-3x)=h(x)$, nên $F(s)=h(-s/3)$ và
\[u=y^3+h(x-y/3).\]
Ba công thức cuối thỏa PDE vì được lấy từ các nghiệm tổng quát đã kiểm tra; thế $y=0$ cho đúng $h(x)$.

\textbf{Trường hợp đặc trưng ở (a).} Nếu $b=0$, trên dữ liệu phải có $ah'(x)=f$, tức $a(4x+3)=f$ ở cả khoảng. Với các hằng số $a,f$, điều này chỉ có thể xảy ra khi $a=f=0$. Khi ấy PDE không ràng buộc và có vô số nghiệm cùng vết $h$. Trong các trường hợp còn lại không có nghiệm thỏa dữ liệu. Khi $b\ne0$, và ở (b),(c),(d), đường $y=0$ không đặc trưng, nên dữ liệu xác định duy nhất các nghiệm trên toàn mặt phẳng.

\textbf{Chi tiết thế đa thức dữ liệu vào bốn nghiệm ở (e).} Với $h(s)=2s^2+3s$, các công thức lần lượt là
\[\begin{aligned}
u_a&=2(x-ay/b)^2+3(x-ay/b)+fy/b &&(b\ne0),\\
u_b&=\tfrac12+e^{-y}[2(x-y/2)^2+3(x-y/2)-\tfrac12],\\
u_c&=\frac{x^2}{6}+\frac{e^x}{3}+\frac{11}{6}(x+3y/4)^2+3(x+3y/4)-\frac{e^{x+3y/4}}3,\\
u_d&=y^3+2(x-y/3)^2+3(x-y/3).
\end{aligned}\]
Ở (b), đặt $z=x-y/2,Q=h(z)-1/2$; $u_x=e^{-y}h'(z)$, $u_y=-e^{-y}Q-e^{-y}h'(z)/2$. Vậy $u_x+2u_y+2u=1$. Ở (c), đặt $z=x+3y/4$, $u_x=P'(x)+h'(z)-P'(z)$, $u_y=3[h'(z)-P'(z)]/4$; thế vào $3u_x-4u_y$ còn $3P'(x)=x+e^x$. Ở (d), $u_x=h'(x-y/3)$, $u_y=3y^2-h'(x-y/3)/3$, cho $u_x+3u_y=9y^2$.

%ID BT_1_3_1_3_04
\textbf{Bước 1 — Giải PDE trước khi dùng dữ liệu.} $dy/dx=2$ cho $\rho=y-2x$. Trên đường, $U'-4U=e^{3x+\rho}$ và $(e^{-4x}U)'=e^{\rho-x}$. Tích phân cho
\[u=-e^{x+y}+e^{4x}F(y-2x).\]

\textbf{Bước 2 — Xác định $F$ từng ý.}
(a) $u(0,y)=y^2$ cho $-e^y+F(y)=y^2$, nên $F(s)=s^2+e^s$.

(b) $u(x,-x)=x$ cho $-1+e^{4x}F(-3x)=x$. Đặt $s=-3x$, suy $F(s)=(1-s/3)e^{4s/3}$.

(c) $u(x,4x)=x^2$ cho $-e^{5x}+e^{4x}F(2x)=x^2$. Đặt $s=2x$, suy $F(s)=s^2e^{-2s}/4+e^{s/2}$.

Trong mỗi ý, thay hàm $F$ vừa tìm vào công thức ở bước 1 để được nghiệm. Chú ý ở (b) cả hạng 1 và hạng $-s/3$ đều nhân $e^{4s/3}$.

(d) Dữ liệu trên parabol cho
\[F(x^2-2x)=H(x),\quad H(x)=e^{-4x}[f(x)+e^{x+x^2}].\]
Vì $x^2-2x=(x-1)^2-1$, hai điểm $x$ và $2-x$ nằm trên cùng đặc trưng. Cần $H(x)=H(2-x)$. Khi tương thích, đặt $F(s)=H(1+\sqrt{s+1})$ trên $s\ge-1$ và kiểm tra độ trơn tại $s=-1$; đối xứng đơn thuần chưa đủ nếu $H$ thiếu trơn. Miền dữ liệu phủ là $y-2x\ge-1$. Muốn nghiệm trên toàn mặt phẳng phải có kéo dài $F$ đủ trơn qua $-1$; trên $s<-1$ kéo dài không được dữ liệu quyết định.

\textbf{Bước 3 — Kiểm tra.} Với $L=\partial_x+2\partial_y-4$, phần $e^{4x}F(\rho)$ có $L=0$, phần $-e^{x+y}$ có $L=e^{x+y}$. Thế lại các đường dữ liệu vào $F$ kiểm tra đúng từng giá trị. Ở (a),(b),(c), chỉ số đặc trưng lần lượt là $y,-3x,2x$, phủ toàn trục nên cho duy nhất.

\textbf{Chi tiết thay $F$ và viết nghiệm cuối của (a),(b),(c).} Trong $u=-e^{x+y}+e^{4x}F(y-2x)$:
\[\begin{aligned}
(a)\quad u&=-e^{x+y}+e^{4x}(y-2x)^2+e^{2x+y},\\
(b)\quad u&=-e^{x+y}+\frac{3+2x-y}{3}e^{4(x+y)/3},\\
(c)\quad u&=-e^{x+y}+\frac{(y-2x)^2}{4}e^{8x-2y}+e^{3x+y/2}.
\end{aligned}\]
Ở (b), $4x+4(y-2x)/3=4(x+y)/3$; ở (c), $4x-2(y-2x)=8x-2y$ và $4x+(y-2x)/2=3x+y/2$. Các bước này giữ nguyên thừa số mũ khi thế $F(s)=s^2e^{-2s}/4+e^{s/2}$.
Thế $x=0$ vào (a): $-e^y+y^2+e^y=y^2$. Thế $y=-x$ vào (b): $-1+(3+3x)/3=x$. Thế $y=4x$ vào $u_c$: $-e^{5x}+x^2+e^{5x}=x^2$.

%ID BT_1_3_1_3_05
\textbf{Bước 1 — Giải dọc đặc trưng.} $u_x+3u_y+u=1$ cho $\rho=y-3x$ và $U'+U=1$. Nhân $e^x$, tích phân được $u=1+e^{-x}F(y-3x)$.

\textbf{Bước 2 — Xét dữ liệu trên $y=3x$.} Trên đường này $\rho=0$, nên $g(x)=1+e^{-x}F(0)$. Do đó điều kiện cần và đủ là $g(x)=1+Ae^{-x}$ với $A$ hằng. Cũng có thể đạo hàm vết để có $g'+g=1$.

\textbf{Bước 3 — Tồn tại và không duy nhất.} Nếu dữ liệu phù hợp, chọn bất kỳ $F$ đủ trơn với $F(0)=A$. Hai lựa chọn $F=A$ và $F=A+s$ cho hai nghiệm khác nhau cùng vết.

\textbf{Ý (b) — Dữ liệu in trong đề.} Với $g=-1+2e^x$, ta có $g'+g=-1+4e^x$, không đồng nhất bằng 1. Vì thế vô nghiệm; không thể viết hai nghiệm cho dữ liệu này. Nếu sửa dữ liệu thành $1+2e^{-x}$, hai nghiệm là $1+2e^{-x}$ và $1+e^{-x}(2+y-3x)$. Đây là minh họa cho đề đã sửa, không phải đáp án của dữ liệu nguyên văn.

%ID BT_1_3_1_3_06
\textbf{Bước 1 — Đặc trưng của $au_x+bu_y=0$.} Giải $X'=a,Y'=b$; tính $(bX-aY)'=0$, nên bất biến $\rho=bx-ay$.

\textbf{Bước 2 — Tích phân ẩn.} $U'=0$ cho $U$ hằng trên mỗi đặc trưng, nhưng hằng có thể khác giữa các đường. Suy $u=F(bx-ay)$ với $(a,b)\ne(0,0)$.

\textbf{Bước 3 — Kiểm tra và ngoại lệ.} $u_x=bF'$, $u_y=-aF'$ nên PDE đúng. Nếu $a=b=0$, phương trình không ràng buộc, mọi hàm đủ trơn là nghiệm; không dùng dạng $F(0)$ để thu hẹp sai họ nghiệm.

%ID BT_1_3_1_3_07
\textbf{Bước 1 — Nghiệm tổng quát.} Với $(a,b)\ne(0,0)$, bài trước cho $u=F(bx-ay)$.

\textbf{Ý (a).} Trên $bx-ay=1$, nghiệm phải bằng hằng $F(1)$. Do đó dữ liệu $h(x)$ phải có cùng giá trị tại mọi điểm thực sự nằm trên đường. Nếu $a\ne0$, $x$ chạy toàn trục và cần $h\equiv H$. Khi đó mọi $F$ với $F(1)=H$ cho nghiệm, nên không duy nhất. Nếu $a=0,b\ne0$, đường là $x=1/b$, dữ liệu chỉ gán giá trị $h(1/b)$; không đòi $h$ hằng ở những $x$ không nằm trên đường.

\textbf{Ý (b).} Đường $x-y=1$ có $y=x-1$, nên $F((b-a)x+a)=h(x)$. Nếu $b\ne a$, đặt $s=(b-a)x+a$:
\[F(s)=h\left(\frac{s-a}{b-a}\right),\quad u=h\left(\frac{bx-ay-a}{b-a}\right).\]
Chỉ số $s$ phủ toàn trục, cho duy nhất. Nếu $b=a\ne0$, chỉ số bằng $a$, cần $h$ hằng; khi tương thích vẫn có vô số $F$ với cùng $F(a)$.

\textbf{Ngoại lệ $a=b=0$.} Ở (a), tập $bx-ay=1$ rỗng, nên không có điều kiện dữ liệu thực sự. Ở (b), chỉ cần một hàm đủ trơn có vết $h$ trên đường; PDE không cho duy nhất. Các kết luận trên giả sử dữ liệu có độ trơn phù hợp.

%ID BT_1_3_1_3_08
\textbf{Bước 1 — Giải PDE.} $u_x=0$ nên $u=F(y)$: nghiệm hằng trên các đường ngang.

\textbf{Bước 2 — Thế dữ liệu.} $u(x,x^2)=x$ đòi $F(x^2)=x$. Với $r>0$, hai điểm $x=\sqrt r$ và $x=-\sqrt r$ cùng cho đối số $r$, nhưng yêu cầu $F(r)=\sqrt r$ và $F(r)=-\sqrt r$.

\textbf{Bước 3 — Kết luận và miền hẹp.} Không có nghiệm trên toàn đường dữ liệu. Nếu chỉ giữ nhánh $x>0$, có $F(y)=\sqrt y$ trên $y>0$. Nếu muốn gồm gốc và nghiệm $C^1$, $F'(y)=1/(2\sqrt y)$ không hữu hạn tại 0, nên cũng không có kéo dài $C^1$ theo công thức này.

%ID BT_1_3_1_3_09
\textbf{Bước 1 — Đặc trưng.} Với $u_x-2u_y=0$, $dy/dx=-2$, nên $\rho=y+2x$ và $u=F(\rho)$.

\textbf{Bước 2 — Dùng dữ liệu.} Trên $y=e^x$, $F(e^x+2x)=(e^x+2x)^2$. Hàm $\phi(x)=e^x+2x$ có $\phi'=e^x+2>0$; khi $x\to-\infty$ thì $\phi\to-\infty$, khi $x\to+\infty$ thì $\phi\to+\infty$. Vì vậy mọi chỉ số đặc trưng được gặp đúng một lần, $F(s)=s^2$.

\textbf{Bước 3 — Kiểm tra.} $u=(y+2x)^2$ có $u_x=4(y+2x)$, $u_y=2(y+2x)$, nên $u_x-2u_y=0$; vết trên $y=e^x$ đúng dữ liệu. Nghiệm duy nhất trên cả mặt phẳng.

%ID BT_1_3_1_3_10
\textbf{Bước 1 — Nghiệm riêng và thuần nhất.} Với $3u_x-2u_y+u=x$, thử $u_p=x-3$ và có $3+(x-3)=x$. Đặt $w=u-u_p$, thì $3w_x-2w_y+w=0$. Bất biến $\rho=2x+3y$, ODE theo $x$ là $W'+W/3=0$, nên
\[u=x-3+e^{-x/3}F(2x+3y).\]

\textbf{Ý (a).} Trên $y=x$, $u=x$ đòi $e^{-x/3}F(5x)=3$. Đặt $s=5x$, $F(s)=3e^{s/15}$, suy
\[u=x-3+3e^{(y-x)/5}.\]
Thế $y=x$ cho $x$; $3u_x-2u_y+u=x$ vì phần mũ có hệ số $-3/5-2/5+1=0$.

\textbf{Ý (b).} Trên $2x+3y=1$, dữ liệu $u=0$ đòi $F(1)=(3-x)e^{x/3}$. Vế trái hằng nhưng đạo hàm vế phải là $-xe^{x/3}/3$, không đồng nhất bằng 0. Do đó không có nghiệm trên toàn đường. Đây là đường đặc trưng, nên phải kiểm tra tính tương thích thay vì khử $F$ như ở (a).

%ID BT_1_3_1_3_11
\textbf{Bước 1 — Nêu giả thiết mô hình.} Hiểu $y$ là tuổi, mỗi cá thể tăng tuổi với tốc độ 1. Mật độ ban đầu $P(y,0)=C(y)\ge0$, tổng dân số $N(t)=\int_0^\infty P(y,t)dy$, sinh ở tuổi 0 cho $P(0,t)=\alpha N(t)$. Nếu $y$ chỉ là tọa độ không gian thì đề còn thiếu luật vận chuyển; giả thiết tuổi là điều cần nói rõ.

\textbf{Bước 2 — Giữ cách ghi chết tuyệt đối $D=k$.} Cân bằng dọc tuổi cho $P_t+P_y=-k$. Đặc trưng từ dữ liệu đầu có $y=\xi+t$ và $P'= -k$, nên trên $y>t$,
\[P(y,t)=C(y-t)-kt.\]
Để mật độ không âm ở mọi tuổi với $t>0$ phải có $C(\xi)\ge kt$ ở mọi $\xi>0$. Khi đó $\int C\ge\int_0^\infty kt\,d\xi=\infty$, trái giả thiết hữu hạn. Vì vậy mô hình chết tuyệt đối không cho mật độ dân số không âm hữu hạn trên mọi tuổi.

\textbf{Bước 3 — Giải mô hình tỉ lệ chết nếu đó là ý định của đề.} Sửa thành $D=kP$, PDE là $P_t+P_y=-kP$. Trên một đặc trưng, $P$ giảm theo $e^{-kt}$. Từ dữ liệu đầu, $P=e^{-kt}C(y-t)$ khi $y\ge t$. Từ biên tuổi 0 ở thời điểm $t-y$, $P=e^{-ky}\alpha N(t-y)$ khi $y<t$.

\textbf{Bước 4 — Tính tổng dân số.} Thế hai nhánh vào tích phân:
\[N(t)=e^{-kt}N_0+\alpha\int_0^t e^{-k(t-s)}N(s)ds,\qquad N_0=\int_0^\infty C.\]
Nhân $e^{kt}$, đạo hàm được $(e^{kt}N)'=\alpha e^{kt}N$, nên $N=N_0e^{(\alpha-k)t}$. Suy
\[P(y,t)=\begin{cases}e^{-kt}C(y-t),&y\ge t,\\ \alpha N_0e^{(\alpha-k)t-\alpha y},&0\le y<t.\end{cases}\]

\textbf{Bước 5 — Kiểm tra và ghép.} Hai nhánh thỏa $P_t+P_y=-kP$; tại $t=0$ đúng $C$, tại $y=0$ đúng $\alpha N(t)$. Muốn liên tục qua $y=t$ cần $C(0)=\alpha N_0$; muốn $C^1$ qua đó cần thêm $C'(0)=-\alpha^2N_0$. Nếu không tương thích, công thức mô tả nghiệm từng miền/ nghiệm vận chuyển yếu. Đây là kết quả của mô hình đã sửa, không đổi ngầm dữ kiện gốc.

%ID VD_1_3_2_2_01
\textbf{Bước 1 — Tìm bất biến của trường hệ số.} Với $Au_x+Bu_y+Cu=R$, giải đặc trưng $X'=A,Y'=B$ để tìm $\rho$ thỏa $A\rho_x+B\rho_y=0$.

\textbf{Bước 2 — Chọn tọa độ còn lại.} Đặt $w=\rho(x,y)$, $z=y$, $u=v(w,z)$. Jacobian là $\rho_x$; chỉ dùng phép đổi ở nơi $\rho_x\ne0$.

\textbf{Bước 3 — Tính đạo hàm và thế.}
\[u_x=v_w\rho_x,\qquad u_y=v_w\rho_y+v_z.\]
Do đó $Au_x+Bu_y=(A\rho_x+B\rho_y)v_w+Bv_z=Bv_z$. PDE thành $Bv_z+Cv=R$, với mọi hệ số được viết theo tọa độ mới. Nếu $B\ne0$, chia và giải ODE tuyến tính theo $z$, xem $w$ là tham số.

\textbf{Bước 4 — Ngoại lệ và dữ liệu.} Nếu $\rho_x=0$, có thể chọn $z=x$ khi $\rho_y\ne0$, hoặc tìm tọa độ khác. Sau giải ODE, đổi về $x,y$, thế đường dữ liệu để tìm hàm tùy ý và kiểm tra PDE. Không dùng phép đổi suy biến để suy tính duy nhất.

%ID VD_1_3_2_2_03
\textbf{Ý (a) — $u_x+x^2u_y=0$.} Giải $dy/dx=x^2$ cho $\rho=y-x^3/3$, rồi $U'=0$ cho $u=F(y-x^3/3)$. Kiểm tra $u_x=-x^2F'$, $u_y=F'$, hai hạng triệt tiêu.

\textbf{Ý (b) — $e^{x^2}u_x+xu_y=0$.} Vì hệ số $e^{x^2}>0$, giải $dy/dx=xe^{-x^2}$, được $\rho=y+e^{-x^2}/2$. Do $U'=0$, $u=F(\rho)$. Kiểm tra $u_x=-xe^{-x^2}F'$, $u_y=F'$, nên PDE bằng 0. Bất biến có thể nhân 2 và đổi tên $F$ để được dạng ở ví dụ 2.

\textbf{Ý (c) — $u_x+\sin x\,u_y=0$.} $dy/dx=\sin x$ cho $y=-\cos x+\rho$. Dọc đường $U'=0$, nên $u=F(y+\cos x)$. Đạo hàm $u_x=-\sin x F'$, $u_y=F'$ kiểm tra PDE.

\textbf{Ý (d) — $xu_x+yu_y=0$.} Giải hệ $X'=X,Y'=Y$: $X=x_0e^s,Y=y_0e^s$. Đặc trưng là các tia, không phải toàn đường thẳng qua gốc; chiều co giãn không đổi dấu của tọa độ. Trên mỗi nửa mặt phẳng $x\ne0$, $Y/X$ không đổi, nên $u=F(y/x)$; trên các miền $y\ne0$, dùng $u=G(x/y)$ để phủ các trục. Các hàm ở các miền cần ghép đủ trơn.

\textbf{Tại gốc.} Nếu yêu cầu nghiệm liên tục trên cả mặt phẳng, với mỗi $(x,y)$ ta có $u(x,y)=u(e^{-s}x,e^{-s}y)$. Cho $s\to\infty$, vế phải tiến tới $u(0,0)$, nên mọi điểm có cùng giá trị. Vì vậy nghiệm toàn mặt phẳng liên tục phải hằng.

%ID BT_1_3_2_3_01
\textbf{Ý (a).} Với $u_x+\sin x\,u_y=0$, tích phân $dy/dx=\sin x$ cho $\rho=y+\cos x$. $U'=0$ nên $u=F(\rho)$. Kiểm tra $u_x+\sin x\,u_y=-\sin xF'+\sin xF'=0$.

\textbf{Ý (b).} Với $xu_x+yu_y=0$, $X=x_0e^s,Y=y_0e^s$, nên $u$ hằng trên tia. Trên từng miền $x>0$ hoặc $x<0$, $u=F(y/x)$; dùng tọa độ $x/y$ trên miền $y\ne0$. Nếu cần liên tục trên toàn mặt phẳng, co giãn mỗi điểm về gốc cho $u(x,y)=u(0,0)$, nên chỉ có nghiệm hằng.

\textbf{Ý (c).} Với $xu_x+2yu_y=0$ trong $x,y>0$, $X'=X,Y'=2Y$. Tính $(Y/X^2)'=0$, nên $\rho=y/x^2$. $U'=0$ cho $u=F(y/x^2)$. Kiểm tra $u_x=-2yF'/x^3$, $u_y=F'/x^2$, hai hạng của PDE triệt tiêu.

\textbf{Ý (d).} Với $yu_x-4xu_y=2xy$, trường $V=y\partial_x-4x\partial_y$ có bất biến $r=y^2+4x^2$ vì $Vr=2y(-4x)+8xy=0$. Chọn nghiệm riêng $P=-y^2/4$, có $VP=2xy$. Đặt $w=u-P$, $Vw=0$. Các mức $r>0$ là ellipse liên thông và đặc trưng kín, nên
\[u=-y^2/4+F(y^2+4x^2).\]
Đạo hàm của phần $F$ cho $V=0$; phần riêng cho đúng nguồn. Để có nghiệm toàn mặt phẳng, tổ hợp xuyên tâm ellipse phải đủ trơn tại gốc; $F\in C^1([0,\infty))$ là một điều kiện đủ, nhưng không phải mô tả duy nhất của lớp trơn ở gốc.

\textbf{Chi tiết hệ đặc trưng ở (d).} Giải $X'=Y,Y'=-4X$: $X''=-4X$, nên $X=A\cos2s+B\sin2s$, $Y=-2A\sin2s+2B\cos2s$. Vì $4X^2+Y^2=4(A^2+B^2)$, hệ đi hết một ellipse sau chu kỳ $\pi$. Với $U'=2XY$, có $(Y^2/4)'=YY'/2=-2XY$, do đó $(U+Y^2/4)'=0$. Hằng số trên ellipse là $F(4X^2+Y^2)$, dẫn tới $u=-y^2/4+F(4x^2+y^2)$.

%ID BT_1_3_2_3_02
\textbf{Bước 1 — Đặc trưng.} $-yu_x+xu_y=0$ có $X'=-Y,Y'=X$. Khi đó $X^2+Y^2$ không đổi, đặc trưng là các vòng tròn tâm gốc và $U'=0$.

\textbf{Bước 2 — So sánh dữ liệu trên cùng vòng.} Vòng bán kính $r>0$ gặp trục $x$ ở $(r,0)$ và $(-r,0)$. Nghiệm phải có hai giá trị bằng nhau, trong khi dữ liệu $u(x,0)=3x$ cho $3r$ và $-3r$.

\textbf{Bước 3 — Kết luận.} Không có nghiệm trên toàn mặt phẳng. Ngay trên trục cũng thấy điều kiện cần $f(r)=f(-r)$ bị vi phạm. Nếu chỉ xét miền cục bộ quanh một phần đường dữ liệu, phải nói rõ miền; kết luận toàn cục không tự áp cho mọi miền nhỏ.

\textbf{Chi tiết tích phân hệ quay.} Từ $X'=-Y,Y'=X$, lấy đạo hàm dòng đầu được $X''+X=0$. Với điểm đầu $(x_0,y_0)$, $X=x_0\cos s-y_0\sin s$, $Y=x_0\sin s+y_0\cos s$. Tính $X^2+Y^2=x_0^2+y_0^2$ và $U'=0$ cho $U=u_0$. Sau góc $\pi$, $(r,0)$ đi tới $(-r,0)$ mà $U$ không đổi; đây là lý do hai giá trị dữ liệu phải bằng nhau.

%ID BT_1_3_2_3_03
\textbf{Bước 1 — Hình dạng nghiệm.} Như bài trước, nghiệm hằng trên mỗi vòng tròn nên $u=H(r)$, $r=\sqrt{x^2+y^2}$, ngoài gốc.

\textbf{Bước 2 — Điều kiện dữ liệu.} Trên trục $x$, $H(|x|)=f(x)$, nên cần $f$ chẵn. Nếu $f\in C^1(\mathbb R)$ chẵn thì $f'(0)=0$, và đặt $H(r)=f(r)$, $r\ge0$.

\textbf{Bước 3 — Kiểm tra độ trơn và PDE.} Với $r>0$, $u_x=f'(r)x/r$, $u_y=f'(r)y/r$, nên $-yu_x+xu_y=0$. Do $f'(r)\to0$, gradient tiến về 0 tại gốc; đạo hàm tại gốc cũng bằng 0, nên công thức cho nghiệm $C^1$ trên toàn mặt phẳng.

\textbf{Bước 4 — Duy nhất.} Mọi bán kính gặp dữ liệu nên $H$ được xác định hết. Ví dụ $f=x^2$ cho $u=x^2+y^2$; $f$ không chẵn cho vô nghiệm. Với miền toàn mặt phẳng của đề, không có trường hợp dữ liệu tương thích mà không duy nhất; điều đó chỉ có thể xảy ra trên miền hoặc dữ liệu bị thu hẹp.

\textbf{Chi tiết từ hệ đặc trưng đến nghiệm xuyên tâm.} Hệ $X'=-Y,Y'=X,U'=0$ có nghiệm $X=x_0\cos s-y_0\sin s$, $Y=x_0\sin s+y_0\cos s$, $U=u_0$. Mỗi quỹ đạo $r>0$ đi hết vòng $X^2+Y^2=r^2$, nên giá trị chỉ có thể phụ thuộc $r$. Ở gốc đặt $u(0,0)=f(0)$. Vì $f$ chẵn và khả vi, $f'(0)=-f'(0)$, suy $f'(0)=0$. Hơn nữa $|u(x,y)-u(0,0)|/r=|f(r)-f(0)|/r\to0$, nên $u$ khả vi ở gốc với gradient $0$; các công thức gradient ngoài gốc tiến về cùng giá trị ấy.

%ID BT_1_3_2_3_04
\textbf{Bước 1 — Giữ đúng hệ số của đề.} PDE là $xu_x-yu_y+yu=0$, có hệ số của $u$ là $y$. Hệ đặc trưng $X'=X,Y'=-Y$ cho $r=XY$ không đổi.

\textbf{Bước 2 — Giải ẩn.} $U'=-YU$, còn $Y'=-Y$, nên $(e^{-Y}U)'=e^{-Y}(U'-Y'U)=0$. Do đó
\[u(x,y)=e^yF(xy).\]
Trên mỗi miền $x>0$ hoặc $x<0$, có thể chọn tọa độ $(r,\log|x|)$, nên công thức mô tả nghiệm trên các phần đặc trưng đó. Khi muốn qua trục, ghép đủ trơn.

\textbf{Bước 3 — Kiểm tra.} $u_x=e^y yF'$, $u_y=e^y(F+xF')$, nên $xu_x-yu_y+yu=e^y[xyF'-yF-xyF'+yF]=0$. Trên $x=0,y\ne0$, PDE cho $u_y=u$, phù hợp $u(0,y)=e^yF(0)$; trên $y=0$, PDE cho $xu_x=0$.

%ID BT_1_3_2_3_05
\textbf{Bước 1 — Đặc trưng và ODE.} Với $xu_x-yu_y+yu=y^2$, bất biến vẫn là $r=xy$. Dọc đặc trưng $U'+YU=Y^2$, $Y'=-Y$.

\textbf{Bước 2 — Tích phân.} Nhân $e^{-Y}$:
\[(e^{-Y}U)'=Y^2e^{-Y}.\]
Vì $dY=-Yds$, tích phân vế phải thành $-\int Ye^{-Y}dY=(Y+1)e^{-Y}$. Suy
\[u=y+1+e^yF(xy).\]
Nghiệm riêng $y+1$ giúp thấy công thức vẫn hợp lệ ở $y=0$, dù bước đổi tham số theo $Y$ cần xét riêng nơi $Y=0$.

\textbf{Bước 3 — Kiểm tra.} $x(y+1)_x-y(y+1)_y+y(y+1)=-y+y^2+y=y^2$; phần $e^yF(xy)$ thỏa thuần nhất theo bài 4. Do đó tổng đúng PDE.

\textbf{Chi tiết tích phân từng phần.} Khi $Y\ne0$, $ds=-dY/Y$, nên $\int Y^2e^{-Y}ds=-\int Ye^{-Y}dY$. Chọn $a=Y,db=e^{-Y}dY$, được $\int Ye^{-Y}dY=-Ye^{-Y}-e^{-Y}=-(Y+1)e^{-Y}$. Vậy $e^{-Y}U=(Y+1)e^{-Y}+F(XY)$, nhân lại $e^Y$ cho công thức đã nêu. Với $u_h=e^yF(xy)$, $u_{h,x}=ye^yF'$, $u_{h,y}=e^yF+xe^yF'$ và $xu_{h,x}-yu_{h,y}+yu_h=0$ bằng triệt tiêu từng hạng.

%ID BT_1_3_2_3_06
\textbf{Bước 1 — Dùng hai nghiệm tổng quát đã kiểm tra.} Bài 4: $u=e^yF(xy)$. Bài 5: $u=y+1+e^yF(xy)$. Xét dữ liệu cho từng phương trình.

\textbf{Ý (a) — $u(x,1/x)=x$, $x>0$.} Trên đường này $xy=1$. Với bài 4 phải có $F(1)=xe^{-1/x}$; đạo hàm vế phải bằng $e^{-1/x}(1+1/x)>0$, nên không thể hằng. Với bài 5 phải có $F(1)=e^{-1/x}(x-1-1/x)$; đạo hàm bằng $e^{-1/x}(1+1/x-1/x^3)$, không đồng nhất bằng 0. Cả hai đều vô nghiệm trên toàn đường dữ liệu.

\textbf{Ý (b) — $u(1,y)=y^2$.} Với bài 4, $F(y)=y^2e^{-y}$, nên
\[u=x^2y^2e^{y-xy}.\]
Với bài 5, $F(y)=e^{-y}(y^2-y-1)$, nên
\[u=y+1+e^{y-xy}(x^2y^2-xy-1).\]
Thế $x=1$ cho $y^2$ ở cả hai. Trong miền $x>0$, mọi đặc trưng gặp $x=1$ đúng một lần nên duy nhất ở miền đó; dữ liệu không tự xác định các đặc trưng trong $x<0$.

\textbf{Ý (c) — $u(x,0)=x^4$.} Với cả hai PDE, đặt $y=0$ trong phương trình gốc cho $xu_x(x,0)=0$. Dữ liệu có $u_x(x,0)=4x^3$, suy $4x^4=0$, mâu thuẫn trên mọi khoảng có $x\ne0$. Vì vậy vô nghiệm. Dùng phương trình gốc ở đường suy biến tránh chia cho $y$.

\textbf{Chi tiết hai công thức dùng ở bước 1.} Đặc trưng $X'=X,Y'=-Y$ cho $XY=\rho$. Với phương trình thuần nhất, $(e^{-Y}U)'=0$ nên $u=e^yF(\rho)$. Với phương trình có nguồn $y^2$, $(e^{-Y}U)'=Y^2e^{-Y}$; đổi $ds=-dY/Y$ và tích phân được $(Y+1)e^{-Y}+F(\rho)$, tức $u=y+1+e^yF(\rho)$. Ở (b), đường $x=1$ cho $\rho=y=s$, nên thay $s=xy$ vào $F(s)=s^2e^{-s}$ hoặc $F(s)=(s^2-s-1)e^{-s}$; đây là phép thế tạo ra hai nghiệm đã viết.

%ID BT_1_3_2_3_07
\textbf{Bước 1 — Bất biến.} Với $-2yu_x+xu_y=0$, $X'=-2Y,Y'=X$. Tính $(X^2+2Y^2)'=2X(-2Y)+4YX=0$.

\textbf{Bước 2 — Giá trị trên ellipse.} $U'=0$, nên nghiệm hằng trên mỗi ellipse $x^2+2y^2=r>0$. Do vòng ellipse liên thông, $u=F(x^2+2y^2)$ ngoài gốc.

\textbf{Bước 3 — Kiểm tra và độ trơn.} $u_x=2xF'$, $u_y=4yF'$, nên $-2yu_x+xu_y=0$. Để mô tả đúng lớp $C^1$ ở gốc, cũng có thể viết $u=H(\sqrt{x^2+2y^2})$ với $H\in C^1$, $H'(0)=0$. Yêu cầu $F'$ hữu hạn tại 0 là điều kiện đủ, nhưng có thể loại một số nghiệm $C^1$ hợp lệ.

\textbf{Chi tiết giải hệ trên ellipse.} Từ $X'=-2Y,Y'=X$, suy $X''=-2X$. Với $\omega=\sqrt2$,
\[X=x_0\cos\omega s-\omega y_0\sin\omega s,\quad
Y=y_0\cos\omega s+\frac{x_0}{\omega}\sin\omega s.\]
Thế vào $X^2+2Y^2$ cho $x_0^2+2y_0^2$. Chu kỳ $2\pi/\sqrt2$ đi hết ellipse; $U'=0$ cho hằng số trên toàn ellipse, thay vì chỉ trên một cung.

%ID BT_1_3_2_3_08
\textbf{Ý (a).} $yu_x-xu_y=u$ cho $X'=Y,Y'=-X,U'=U$. Đường đi là vòng tròn chu kỳ $2\pi$, nhưng $U(s+2\pi)=e^{2\pi}U(s)$. Một nghiệm đơn trị phải trở lại cùng giá trị, buộc $U=0$ trên mọi vòng. Dữ liệu $u(x,0)=x$ khác 0, nên vô nghiệm toàn cục. Cục bộ trong một góc không phủ vòng kín có thể có nghiệm, nhưng không kéo dài thành nghiệm đơn trị quanh vòng.

\textbf{Ý (b).} $u_t+xtu_x=0$. Dùng $t$: $X'=tX$, $U'=0$, $X(0)=\xi$, $U(0)=\sin\xi$. Tích phân $X=\xi e^{t^2/2}$, nên $\xi=xe^{-t^2/2}$ và
\[u=\sin(xe^{-t^2/2}).\]
Đặt $z=xe^{-t^2/2}$, $u_t=-tz\cos z$, $xtu_x=tz\cos z$, tổng bằng 0; tại $t=0$ đúng $\sin x$.

\textbf{Ý (c) — Giữ hệ số $-2xt$ trong đề.} $xu_t-2xtu_x=2tu$, $u(x,0)=x^2$. Trên $x\ne0$, chia cho $x$: $u_t-2tu_x=(2t/x)u$. Đặc trưng $X'= -2t$ cho $X=\xi-t^2$, nên $\xi=x+t^2$. Trên đường,
\[U'=\frac{2t}{\xi-t^2}U,\quad \frac{d}{dt}[(\xi-t^2)U]=0.\]
Ở $t=0$, tích này bằng $\xi^3$, do đó
\[u=\frac{(x+t^2)^3}{x}.\]
Đối với miền từ dữ liệu $x>0$, công thức hợp lệ khi $x>0$; từ dữ liệu $x<0$, cần $\xi=x+t^2<0$. Đặt $z=x+t^2$, $u_t=6tz^2/x$, $u_x=3z^2/x-z^3/x^2$, nên $xu_t-2xtu_x=2tz^3/x=2tu$. Trên $x=0,t>0$ phương trình gốc buộc $u=0$, nhưng công thức phía $x>0$ phân kỳ khi $x\downarrow0$. Vì thế không có nghiệm $C^1$ trên toàn mặt phẳng thời gian dương thỏa toàn bộ dữ liệu.

\textbf{Ý (d).} $yu_x+xu_y=0$ có $r=x^2-y^2$ bất biến và $U'=0$. Trên $x=0$, $r=-y^2$; dữ liệu cho $F(r)=e^r$ với $r\le0$. Vì thế trên miền các đặc trưng từ trục $y$ phủ, có $u=e^{x^2-y^2}$. Công thức cũng là một nghiệm toàn mặt phẳng, nhưng miền $r>0$ chưa được dữ liệu xác định. Chẳng hạn thêm $\psi(r)$ với $\psi(r)=0$ khi $r\le0$, $\psi(r)=e^{-1/r}$ khi $r>0$, vẫn cho nghiệm trơn cùng dữ liệu và khác trên $r>0$.

\textbf{Ý (e).} Với $V=y\partial_x+x\partial_y$, $Vr=0$ và $V(y^2/2)=xy$. Do đó $u=y^2/2+H(r)$. Trên $y=0$, $H(r)=e^{-r}$ với $r\ge0$. Trên $x=0$, $r=-y^2$, nên $H(r)=e^r+r/2$ với $r\le0$. Hai giá trị bằng nhau tại 0, nhưng $H'_+(0)=-1$, $H'_-(0)=3/2$. Trên đường $x=y>0$, $\nabla r\ne0$, nên gradient của $u$ không ghép liên tục. Kết luận: công thức ghép là nghiệm liên tục từng miền, nhưng không có nghiệm $C^1$ trên toàn góc phần tư thỏa hai vết đã cho.

\textbf{Chi tiết hệ hyperbola của (d),(e).} Với $X'=Y,Y'=X$, đặt $S=X+Y,T=X-Y$. Khi đó $S'=S,T'=-T$, nên $S=S_0e^s,T=T_0e^{-s}$ và $ST=X^2-Y^2=\rho$. Ở (d), $U'=0$ nên $U$ hằng trên từng thành phần hyperbola; trên trục $x=0$ có $\rho=-y^2$ và dữ liệu $e^{-y^2}=e^\rho$, quyết định các thành phần từ trục ấy. Ở (e), $U'=XY$ và $(Y^2/2)'=YY'=XY$, nên $(U-Y^2/2)'=0$. Bởi vậy $u=y^2/2+H(\rho)$ trước khi dùng hai vết. Đạo hàm hai biểu thức của $H$ tại $0$: $(e^{-r})'|_0=-1$, $(e^r+r/2)'|_0=1+1/2=3/2$, giải thích sự không ghép $C^1$.

%ID VD_1_4_2_01
\textbf{Ý (a) — $u_t+tu_x=u^2$.} Bước 1: $X'=t$ cho $\rho=x-t^2/2$. Bước 2: $U'=U^2$. Nếu $U\ne0$, $(1/U)'=-1$, nên $u=1/[F(\rho)-t]$, ở nơi mẫu khác 0. Phải thêm nghiệm $u=0$ bị mất khi chia. Một dạng từ vết $f(\rho)$ tại $t=0$ giữ cả điểm $f=0$ là $u=f(\rho)/[1-tf(\rho)]$. Dọc đặc trưng công thức có đạo hàm bằng $U^2$, nên kiểm tra được PDE.

\textbf{Ý (b) — $yu_x+xu_y=u^2+u$.} $V=y\partial_x+x\partial_y$ có bất biến $r=x^2-y^2$. Trên $x+y\ne0$, đặt $s=\log|x+y|$; $Vs=1$. ODE là $U'=U(U+1)$. Nếu $U\ne0,-1$, tích phân $dU/[U(U+1)]=ds$ cho $\log|U/(U+1)|=s+\text{hằng}$. Suy
\[u=\frac{A(r)e^s}{1-A(r)e^s}.\]
$u=0$ nằm trong họ khi $A=0$; thêm $u=-1$. Công thức trên mỗi biểu đồ có mẫu khác 0; trên $x+y=0$ dùng biểu đồ khác nếu trường không triệt tiêu, tại gốc PDE cho $u=0$ hoặc $-1$. Ghép đủ trơn trước khi kết luận toàn cục.

\textbf{Ý (c) — $u_x+xyu_y=1/u$.} $dy/dx=xy$ cho $r=ye^{-x^2/2}$. Dọc đường, $U'=1/U$, nên $(U^2)'=2$, $U^2=2x+A(r)$. Suy $u=\pm\sqrt{2x+A(ye^{-x^2/2})}$. Chọn miền trong căn dương và dấu liên tục. Kiểm tra bằng cách đạo hàm $u^2$: $(\partial_x+xy\partial_y)u^2=2$, suy PDE vì $u\ne0$. Không thêm $u=0$ vì phương trình gốc chứa $1/u$.

\textbf{Chi tiết phân thức ở (b).} $\frac1{U(U+1)}=\frac1U-\frac1{U+1}$, nên $\int dU/[U(U+1)]=\log|U|-\log|U+1|$. Sau khi tích phân theo $s$, viết $U/(U+1)=A(r)e^s$ với dấu được gộp vào $A$. Nhân chéo: $U=Ae^sU+Ae^s$, tức $(1-Ae^s)U=Ae^s$; chia khi mẫu khác $0$ được công thức. Hai giá trị cân bằng $0,-1$ phải được xét trước phép chia.

%ID VD_1_4_2_02
\textbf{Ý (a).} $u_t+uu_x=0$, $u(x,0)=3x$. $U'=0$ cho $U=3\xi$; $X'=U$ cho $X=\xi(1+3t)$. Khử $\xi=x/(1+3t)$:
\[u=3x/(1+3t).\]
$X_\xi=1+3t>0$ khi $t\ge0$, nên nghiệm cổ điển trên toàn miền thời gian dương. $u_t=-9x/(1+3t)^2$, $uu_x=9x/(1+3t)^2$.

\textbf{Ý (b).} $u_t-u^2u_x=3u$, $u(x,0)=f(x)$. Giải $U'=3U$ trước: $U=e^{3t}f(\xi)$. Sau đó
\[X'= -e^{6t}f(\xi)^2,\quad X=\xi-\frac{e^{6t}-1}{6}f(\xi)^2.\]
Nghiệm là $u=e^{3t}f(\xi)$, với $\xi$ xác định từ $x=X$. Tính $X_\xi=1-(e^{6t}-1)f(\xi)f'(\xi)/3$; chỉ khử được cục bộ ở nơi khác 0. Giữ nhãn và phương trình khử là cách giải phù hợp khi $f$ tùy ý.

\textbf{Ý (c).} $u_t-t^2uu_x=-u$. $U=e^{-t}f(\xi)$; $X'= -t^2e^{-t}f(\xi)$. Dùng tích phân từng phần hai lần:
\[\int_0^t s^2e^{-s}ds=2-(t^2+2t+2)e^{-t}=:J(t).\]
Suy $X=\xi-J(t)f(\xi)$, $u=e^{-t}f(\xi)$. Miền khử cần $1-J(t)f'(\xi)\ne0$; tại $t=0$, $J=0$, cho đúng dữ liệu.

\textbf{Ý (d).} $u_t+t^2uu_x=5$, $u(x,0)=x$. $U'=5$ cho $U=\xi+5t$. Tích phân $X'=t^2(\xi+5t)$:
\[X=\xi(1+t^3/3)+5t^4/4,\quad \xi=\frac{x-5t^4/4}{1+t^3/3}.\]
Do đó $u=(x-5t^4/4)/(1+t^3/3)+5t$. Với $t\ge0$, $X_\xi=1+t^3/3>0$. Trên mỗi nhánh trên, $dU/dt=u_t+X'u_x$ bằng đúng nguồn của PDE và dữ liệu tại 0 được giữ, nên công thức tham số thỏa bài.

\textbf{Chi tiết tích phân ở (c).} Đặt $J(t)=\int_0^t s^2e^{-s}ds$. Lần thứ nhất:
\[J=-t^2e^{-t}+2\int_0^t se^{-s}ds.\]
Lần thứ hai:
\[\int_0^t se^{-s}ds=-te^{-t}+\int_0^t e^{-s}ds=-te^{-t}+1-e^{-t}.\]
Thế lại được $J=2-(t^2+2t+2)e^{-t}$, rồi $X-\xi=-f(\xi)J(t)$. Ở (b), tích phân tương tự trực tiếp: $X-\xi=-f(\xi)^2\int_0^te^{6s}ds=-f(\xi)^2(e^{6t}-1)/6$.

%ID BT_1_4_3_01
\textbf{Bước 1 — Chọn phương pháp từng dữ liệu.} Trong vùng trơn dùng $U'=0$, $X'=F'(U)$; dữ liệu nhảy dùng sốc/quạt. Các thông lượng lần lượt là $u^2/2$ cho (a),(b), $u^3/3$ cho (c), $2u^2$ cho (d).

\textbf{Ý (a).} Dữ liệu nhảy từ 2 lên 12 tại $x=1$. Với Burgers, tốc độ trái 2, phải 12, nên hai họ tách xa và tạo quạt. Trong quạt, $F'(u)=u=(x-1)/t$. Do đó
\[u=\begin{cases}2,&x\le1+2t,\\(x-1)/t,&1+2t<x<1+12t,\\12,&x\ge1+12t.\end{cases}\]
Các mép cho giá trị 2,12; trong quạt $u_t=-(x-1)/t^2$, $uu_x=(x-1)/t^2$, nên PDE đúng.

\textbf{Ý (b).} Dữ liệu $u_0(\xi)=-\xi$ trơn. $U=-\xi$, $X=\xi(1-t)$. Khử $\xi=x/(1-t)$ cho $u=-x/(1-t)$ khi $0\le t<1$. $X_\xi=1-t$ triệt tiêu ở $t=1$, mọi đặc trưng hội tụ ở $x=0$. Vì vậy không kéo dài thành nghiệm cổ điển qua thời điểm đó.

\textbf{Ý (c).} Trạng thái 4 trái, 3 phải ở $x=0$; $F'=u^2$ có tốc độ 16 và 9, hội tụ. Rankine–Hugoniot:
\[\sigma=\frac{4^3/3-3^3/3}{4-3}=37/3.\]
Suy $s=37t/3$, $u=4$ khi $x<s$, $u=3$ khi $x>s$. Entropy: $9<37/3<16$.

\textbf{Ý (d).} Trạng thái 3 trái, 2 phải tại $x=1$, $F=2u^2$, tốc độ $F'=4u$ là 12 và 8. $\sigma=[2\cdot9-2\cdot4]/(3-2)=10$, nên $s=1+10t$, $u=3$ trái sốc, 2 phải sốc; $8<10<12$. Các trạng thái hằng thỏa PDE trong mỗi miền, đường sốc thỏa bảo toàn, và giới hạn $t\downarrow0$ khớp dữ liệu.

%ID BT_1_4_3_02
\textbf{Bước 1 — Các sốc trước tương tác.} Burgers có dữ liệu lần lượt 2,1,0 với bước nhảy tại 0 và 1. Cả hai bước đều giảm nên tạo sốc. Tốc độ là trung bình hai trạng thái: $\sigma_1=(2+1)/2=3/2$, $\sigma_2=(1+0)/2=1/2$. Vì vậy $s_1=3t/2$, $s_2=1+t/2$.

\textbf{Bước 2 — Thời điểm gặp.} Giải $3t/2=1+t/2$ cho $t_*=1$, $x_*=3/2$. Khi $0<t<1$,
\[u=\begin{cases}2,&x<s_1,\\1,&s_1<x<s_2,\\0,&x>s_2.\end{cases}\]

\textbf{Bước 3 — Tính lại sốc sau gặp.} Trạng thái giữa biến mất; sốc nối 2 với 0 có tốc độ 1. Từ $s(1)=3/2$, $s=3/2+(t-1)=t+1/2$. Với $t\ge1$, $u=2$ khi $x<s$, $u=0$ khi $x>s$.

\textbf{Bước 4 — Kiểm tra.} Các tốc độ đều nằm giữa tốc độ đặc trưng hai phía: $1<3/2<2$, $0<1/2<1$, $0<1<2$. Vị trí sốc ghép liên tục ở $t=1$ và giới hạn ban đầu đúng các đoạn dữ liệu.

%ID BT_1_5_3_01
\textbf{Bước 1 — Tìm gradient trên dữ liệu.} $F=p^2+pq-5u$, dữ liệu $u(0,s)=4s^2$ cho $q_0=8s$. PDE cho $p_0^2+8sp_0-20s^2=0=(p_0-2s)(p_0+10s)$, nên $p_0=2s$ hoặc $-10s$.

\textbf{Bước 2 — Giải Charpit.} $F_p=2p+q,F_q=p,F_u=-5$, nên $p'=5p,q'=5q,U'=10U$. Đặt $E=e^{5\tau}$:
\[p=p_0E,\quad q=8sE,\quad U=4s^2E^2,\quad X=(2p_0+8s)(E-1)/5,\quad Y=s+p_0(E-1)/5.\]

\textbf{Bước 3 — Khử từng nhánh.} Nếu $p_0=2s$, $X/2+2Y=2sE$, nên $u=(x/2+2y)^2$. Nếu $p_0=-10s$, $5X/2-2Y=-2sE$, nên $u=(5x/2-2y)^2$.

\textbf{Bước 4 — Kiểm tra.} Với $u=L^2$, $L=ax+by$, có $p^2+pq=4a(a+b)L^2$. Hai bộ $(a,b)=(1/2,2),(5/2,-2)$ đều có $4a(a+b)=5$, nên PDE đúng; ở $x=0$ đều bằng $4y^2$. Định thức trên dữ liệu $x=0$ là $-(2p_0+q_0)=\mp12s$, nên mỗi nhánh không đặc trưng ở $s\ne0$; không suy duy nhất ở gốc.

\textbf{Chi tiết tích phân và khử.} $p'=5p$ cho $p=p_0e^{5\tau}$, $q'=5q$ cho $q=8se^{5\tau}$; $U'=p(2p+q)=2(p^2+pq)=10U$ cho $U=4s^2e^{10\tau}$. Vì $\int_0^\tau e^{5r}dr=(E-1)/5$, tích phân $X',Y'$ được các công thức ở bước 2. Với $p_0=2s$: $X=12s(E-1)/5$, $Y=s+2s(E-1)/5$, nên $X/2+2Y=2s+2s(E-1)=2sE$. Với $p_0=-10s$: $X=-12s(E-1)/5$, $Y=s-2s(E-1)$, nên $5X/2-2Y=-2sE$. Bình phương cả hai cho $U$.

%ID BT_1_5_3_02
\textbf{Bước 1 — Dữ liệu trên đường chéo.} $F=p^2-5q^2-u$, $u(s,s)=s^2$ cho $p_0+q_0=2s$. Thế $p_0=2s-q_0$:
\[(2s-q_0)^2-5q_0^2=s^2\ \Longleftrightarrow\ 4q_0^2+4sq_0-3s^2=0.\]
Hai nhánh là $(p_0,q_0)=(7s/2,-3s/2)$ và $(3s/2,s/2)$.

\textbf{Bước 2 — Hệ đặc trưng.} $X'=2p,Y'=-10q,p'=p,q'=q,U'=2U$. Đặt $E=e^\tau$, tích phân:
\[X=s+2p_0(E-1),\quad Y=s-10q_0(E-1),\quad U=s^2E^2.\]

\textbf{Bước 3 — Khử tham số.} Nhánh thứ nhất có $(7X-3Y)/4=sE$, nhánh thứ hai có $(3X+Y)/4=sE$. Vì vậy
\[u=((7x-3y)/4)^2,\qquad u=((3x+y)/4)^2.\]

\textbf{Bước 4 — Kiểm tra.} Với $u=(ax+by)^2$, PDE cần $4(a^2-5b^2)=1$, đúng cho $(7/4,-3/4)$ và $(3/4,1/4)$; $a+b=1$ cho vết $x^2$. $J_0=-10q_0-2p_0=\pm8s$, nên hai nhánh không đặc trưng ở $s\ne0$.

\textbf{Chi tiết nghiệm bậc hai và khử.} $4q_0^2+4sq_0-3s^2=(2q_0-s)(2q_0+3s)$, nên $q_0=s/2,-3s/2$; lấy $p_0=2s-q_0$. Nhánh $(7s/2,-3s/2)$ có $X=s+7s(E-1)$, $Y=s+15s(E-1)$, do đó $7X-3Y=4s+4s(E-1)=4sE$. Nhánh $(3s/2,s/2)$ có $X=s+3s(E-1)$, $Y=s-5s(E-1)$, nên $3X+Y=4sE$. Do $U=s^2E^2$, mỗi tổ hợp chia $4$ rồi bình phương cho nghiệm tương ứng.

%ID BT_1_5_3_03
\textbf{Bước 1 — Giữ hệ số $y$ của đề.} $F=p^2+8qu-8yq^2$. Dữ liệu $(s,s,s^2)$ cho $p_0+q_0=2s$. Đặt $q_0=Q$, $p_0=2s-Q$, ta được
\[(1-8s)Q^2+(8s^2-4s)Q+4s^2=0.\]
Biệt thức là $\Delta=64s^3(s+1)$. Vì $\Delta<0$ trên $-1<s<0$, không có nghiệm thực trên lân cận toàn đường dữ liệu. Với $s>0$ hoặc $s<-1$, lấy các nhánh thực
\[Q=\frac{4s-8s^2\pm8\sqrt{s^3(s+1)}}{2(1-8s)}\quad(s\ne1/8).\]
Ở $s=1/8$, phương trình thành tuyến tính và cho $Q=1/6$; chọn nhánh có giới hạn hữu hạn qua điểm này.

\textbf{Bước 2 — Tính và giải hệ Charpit.} $F_x=0,F_y=-8q^2,F_u=8q,F_p=2p,F_q=8u-16yq$. Do đó
\[X'=2p,\quad Y'=8U-16Yq,\quad p'=-8pq,\quad q'=0.\]
Trên $F=0$, $U'=2p^2+8qU-16Yq^2=-8qU$. Vì $q=Q(s)$ hằng trên đường, đặt $E=e^{-8Q\tau}$:
\[p=p_0E,\quad q=Q,\quad U=s^2E,\quad X=s+\frac{p_0}{4Q}(1-E).\]
Giải $Y'+16QY=8s^2E$ bằng thừa số $e^{16Q\tau}$:
\[Y=sE^2+\frac{s^2}{Q}(E-E^2).\]

\textbf{Bước 3 — Nghiệm và miền.} Nghiệm là $u(X(s,\tau),Y(s,\tau))=s^2E$, với từng nhánh $Q$ trơn, tại nơi ánh xạ khả nghịch. Ở dữ liệu,
\[J_0=8s^2-16sQ-2p_0=8s^2-4s+2(1-8s)Q=\pm\sqrt\Delta.\]
Vì vậy có nghiệm cục bộ trên từng nhánh khi $s>0$ hoặc $s<-1$; $s=0,-1$ suy biến cần xét riêng. Trên các khoảng không đặc trưng này, $Q\ne0$ nên các công thức chia cho $Q$ hợp lệ.

\textbf{Bước 4 — Kiểm tra.} Tại $\tau=0$, $X=Y=s,U=s^2$. Thế các công thức vào $p^2+8qU-8Yq^2$ cho $E^2[p_0^2+8Qs^2-8sQ^2]=0$, theo phương trình tương thích. Hệ Charpit cùng điều kiện tiếp tuyến bảo đảm $p=u_x,q=u_y$ trên miền khả nghịch. Không bỏ hệ số $y$ để dùng công thức của một PDE khác.

\textbf{Chi tiết biệt thức và tích phân $Y$.} Với $a=1-8s,b=8s^2-4s,c=4s^2$,
\[\Delta=b^2-4ac=64s^4-64s^3+16s^2-16s^2+128s^3=64s^3(s+1).\]
Với $Q\ne0$, nhân $Y'+16QY=8s^2e^{-8Q\tau}$ bởi $e^{16Q\tau}$:
\[e^{16Q\tau}Y-s=8s^2\int_0^\tau e^{8Qr}dr=\frac{s^2}{Q}(e^{8Q\tau}-1).\]
Nhân $e^{-16Q\tau}$ được $Y=sE^2+(s^2/Q)(E-E^2)$. Đồng thời $X-s=2p_0\int_0^\tau e^{-8Qr}dr=p_0(1-E)/(4Q)$. Các hằng số đều lấy từ dữ liệu tại $\tau=0$.

%ID BT_1_5_3_04
\textbf{Bước 1 — Tương thích.} $F=p^2+q^2+2xp+2yq-2u$. Dữ liệu $u(s,0)=(1-s^2)/2$ cho $p_0=-s$. Thế PDE: $s^2+q_0^2-2s^2=1-s^2$, suy $q_0^2=1$, $q_0=\pm1$.

\textbf{Bước 2 — Giải Charpit.} $p'=-2p+2p=0$, $q'=-2q+2q=0$. $X'=2p+2X=2(X-s)$, $X(0)=s$, nên $X=s$. $Y'=2q_0+2Y$ cho $Y=q_0(e^{2\tau}-1)$. $U'=2+2q_0Y=2e^{2\tau}$, nên $U=(1-s^2)/2+e^{2\tau}-1$.

\textbf{Bước 3 — Khử và kiểm tra.} $e^{2\tau}-1=q_0Y$, nên
\[u=(1-x^2)/2+y,\qquad u=(1-x^2)/2-y.\]
Với $p=-x,q=\pm1$, vế trái PDE bằng $1-x^2\pm2y=2u$ và vết đúng. $J_0=2q_0\ne0$ nên mỗi nhánh cho nghiệm cục bộ duy nhất khi gradient pháp tuyến đã chọn.

%ID BT_1_5_3_05
\textbf{Bước 1 — Gradient ban đầu.} $F=p^2+q^2+2yp+2xq-2u$. Dữ liệu $u(s,0)=s^2/2$ cho $p_0=s$ và $q_0^2+2sq_0=0$, nên $q_0=0$ hoặc $-2s$.

\textbf{Bước 2 — Tách hệ Charpit.} Tính
\[X'=2p+2Y,\quad Y'=2q+2X,\quad p'=2p-2q,\quad q'=2q-2p.\]
Đặt $R=p+q$, $D=p-q$, $S=X+Y$, $T=X-Y$. Khi đó $R'=0,D'=4D$, $S'=2S+2R$, $T'=-2T+2D$. Với $R_0=s+q_0,D_0=s-q_0,S_0=T_0=s$, tích phân:
\[R=R_0,\quad D=D_0e^{4\tau},\quad S=(s+R_0)e^{2\tau}-R_0,\]
\[T=(s-D_0/3)e^{-2\tau}+(D_0/3)e^{4\tau}.\]
Lấy $X=(S+T)/2$, $Y=(S-T)/2$, $p=(R+D)/2$, $q=(R-D)/2$. Từ PDE,
\[U=(p^2+q^2+2Yp+2Xq)/2.\]
Đây là công thức tham số đã tích phân bằng các hàm mũ thông thường.

\textbf{Bước 3 — Hai nhánh.} Với $q_0=0$, $R_0=D_0=s$, dùng các công thức trên và khử $s,\tau$ ở nơi Jacobian khác 0; không cần ép thành một đa thức. Với $q_0=-2s$, $R_0=-s,D_0=3s$, nên $S=s,T=se^{4\tau}$, và
\[U=(3T^2-S^2)/4,\qquad u=(x^2+y^2)/2-2xy.\]

\textbf{Bước 4 — Kiểm tra và miền.} Nhánh đa thức có $p=x-2y,q=y-2x$; thế cho $p^2+q^2+2yp+2xq=x^2+y^2-4xy=2u$, vết đúng $x^2/2$. Với nhánh tham số, hệ được tích phân từ đúng Charpit và $U$ được khôi phục từ $F=0$; tại 0 đúng toàn bộ dữ liệu. $J_0=2(q_0+s)=\pm2s$, nên mỗi nhánh có thể khử cục bộ ở $s\ne0$; gốc không được định lý này bảo đảm.

\textbf{Chi tiết thừa số tích phân.} $R'=0$ cho $R=R_0$; $D'=4D$ cho $D=D_0e^{4\tau}$. Với $S'-2S=2R_0$, có $(e^{-2\tau}S)'=2R_0e^{-2\tau}$, nên $e^{-2\tau}S=s+R_0(1-e^{-2\tau})$. Với $T'+2T=2D_0e^{4\tau}$, có $(e^{2\tau}T)'=2D_0e^{6\tau}$, nên $e^{2\tau}T=s+D_0(e^{6\tau}-1)/3$. Đây là hai công thức $S,T$ trong bước 2. Ở nhánh $R=-S,D=3T$, $p=(-S+3T)/2,q=(-S-3T)/2$, và
\[p^2+q^2=(S^2+9T^2)/2,\quad 2Yp+2Xq=-S^2-3T^2.\]
Vậy $2U=(-S^2+3T^2)/2$, tức $U=(3T^2-S^2)/4$.

%ID BT_1_5_3_06
\textbf{Bước 1 — Tương thích.} $F=4p^2+q^2+4(p-x)(q-4y)-8u$, dữ liệu $u(s,0)=0$ cho $p_0=0$ và $q_0(q_0-4s)=0$.

\textbf{Bước 2 — Hệ Charpit và tổ hợp bất biến.}
\[X'=8p+4q-16Y,\quad Y'=4p+2q-4X,\quad p'=X',\quad q'=4Y'.\]
Suy $X-p=s$, $q-4Y=q_0$. Thế $p=X-s$, $q=4Y+q_0$ vào hai phương trình đầu:
\[X'=8X-8s+4q_0,\quad Y'=8Y+2q_0-4s.\]
Tích phân với $X(0)=s,Y(0)=0$, đặt $E=e^{8\tau}$:
\[X=s+(q_0/2)(E-1),\quad Y=(q_0/4-s/2)(E-1).\]

\textbf{Bước 3 — Khử.} Với $q_0=0$, $X=s,p=0,q=4Y$, PDE cho $u=2y^2$. Với $q_0=4s$, $X=s+4Y$, nên $s=X-4Y$, $p=4Y,q=4X-12Y$, PDE cho $u=4xy-6y^2$.

\textbf{Bước 4 — Kiểm tra.} Nhánh đầu: $p=0,q=4y$ cho vế trái $16y^2=8u$. Nhánh sau: $p=4y,q=4x-12y$ cho vế trái $32xy-48y^2=8u$. Cả hai bằng 0 ở $y=0$. $J_0=2q_0-4s=\mp4s$ khác 0 khi $s\ne0$; hai nhánh trùng ở gốc không cho suy duy nhất tại đó.

\textbf{Chi tiết tích phân và rút gọn nhánh thứ hai.} Vì $p'=X'$, $p-X=-s$; vì $q'=4Y'$, $q-4Y=q_0$. Giải $X'-8X=-8s+4q_0$ bằng nhân $e^{-8\tau}$ cho $X=s+q_0(E-1)/2$. Tương tự $Y'-8Y=2q_0-4s$ cho $Y=(q_0/4-s/2)(E-1)$. Khi $q_0=4s$, $X-s=2s(E-1)=4Y$, nên $p=4Y,q=4X-12Y$. Thế vào PDE:
\[4p^2+q^2+4(p-X)(q-4Y)
=64Y^2+(4X-12Y)^2-16(X-4Y)^2
=32XY-48Y^2.\]
Chia $8$ cho $U=4XY-6Y^2$.

%ID BT_1_5_3_07
\textbf{Bước 1 — Tương thích và miền log.} $F=pq-1$, $u(s,0)=\log s$ với $s>0$ cho $p_0=1/s$, $q_0=s$.

\textbf{Bước 2 — Giải Charpit.} $p'=q'=0$, $X'=q=s$, $Y'=p=1/s$, $U'=2pq=2$. Do đó $X=s+s\tau$, $Y=\tau/s$, $U=\log s+2\tau$.

\textbf{Bước 3 — Khử.} $\tau=sy$ và $x=s+s^2y$. Nghiệm của phương trình bậc hai liên tục về $s=x$ khi $y=0$ là
\[s=\frac{2x}{1+\sqrt{1+4xy}},\qquad u=\log s+2sy.\]
Chọn miền gần $x>0,y=0$ để $s>0$ và $1+4xy>0$.

\textbf{Bước 4 — Kiểm tra.} Từ $x=s+s^2y$, $s_x=1/(1+2sy)$ và $s_y=-s^2/(1+2sy)$. Đạo hàm $u$ cho $u_x=(1/s+2y)s_x=1/s$, $u_y=(1/s+2y)s_y+2s=s$, nên $pq=1$. Tại $y=0$, $s=x$, $u=\log x$. Mẫu $1+2sy$ cũng là điều kiện khử không suy biến.

\textbf{Chi tiết chọn nghiệm bậc hai.} Khi $y\ne0$, $ys^2+s-x=0$ cho $s=(-1\pm\sqrt{1+4xy})/(2y)$. Nhánh dấu cộng có giới hạn $x$ khi $y\to0$; nhân tử số và mẫu với $1+\sqrt{1+4xy}$ cho $s=2x/(1+\sqrt{1+4xy})$. Nhánh dấu trừ không về dữ liệu $s=x>0$. Sau đó $\tau=sy$ và $U=\log s+2\tau$ cho nghiệm.

%ID BT_1_5_3_08
\textbf{Bước 1 — Dữ liệu trên $y=1$.} $F=pq-u$, $u(s,1)=s$ cho $p_0=1$, $q_0=s$.

\textbf{Bước 2 — Hệ đặc trưng.} $p'=p,q'=q$, $X'=q,Y'=p$, $U'=2U$. Tích phân:
\[p=e^\tau,\quad q=se^\tau,\quad X=se^\tau,\quad Y=e^\tau,\quad U=se^{2\tau}.\]

\textbf{Bước 3 — Khử và thử lại.} $U=XY$, nên $u=xy$. $u_x=y,u_y=x$ cho $pq=xy=u$, và $u(x,1)=x$. Vì dữ liệu chỉ có $0\le s\le1$, miền được đặc trưng phủ cần $y>0$ và $0\le x/y\le1$. Đa thức $xy$ là một kéo dài toàn cục; dữ liệu đoạn không tự xác định duy nhất ở mọi điểm ngoài miền phủ.

%ID BT_1_5_3_09
\textbf{Bước 1 — Tương thích.} $F=pq-2$, $u(s,s)=3s$ cho $p_0+q_0=3$, $p_0q_0=2$. Hai nghiệm của $z^2-3z+2=0$ là 1,2, nên $(p_0,q_0)=(1,2)$ hoặc $(2,1)$.

\textbf{Bước 2 — Tích phân.} $p,q$ hằng; $X=s+q_0\tau$, $Y=s+p_0\tau$, $U=3s+4\tau$. Với nhánh $(1,2)$, $U=X+2Y$; nhánh $(2,1)$ cho $U=2X+Y$.

\textbf{Bước 3 — Nghiệm và miền.} $u=x+2y$ hoặc $u=2x+y$. Gradient có tích 2, vết bằng $3x$, nên cả hai đúng. $J_0=p_0-q_0=\mp1$ khác 0. Từ đoạn $0\le s\le1$, nhánh đầu cần $s=2y-x\in[0,1]$, nhánh sau cần $s=2x-y\in[0,1]$. Các công thức affine kéo dài được, nhưng kết luận duy nhất cục bộ áp cho từng nhánh đã chọn.

%ID BT_1_5_3_10
\textbf{Bước 1 — Gradient đầu.} $F=p^3-q$, $u(s,0)=2s^{3/2}$, $0\le s\le1$, cho $p_0=3\sqrt s$, $q_0=p_0^3=27s^{3/2}$.

\textbf{Bước 2 — Charpit.} $p,q$ hằng, $X'=3p^2=27s$, $Y'=-1$, $U'=3p^3-q=54s^{3/2}$. Suy $X=s+27s\tau$, $Y=-\tau$, $U=2s^{3/2}(1+27\tau)$.

\textbf{Bước 3 — Khử.} $\tau=-y$, $s=x/(1-27y)$, nên
\[u=2x^{3/2}/\sqrt{1-27y}.\]
Chọn miền gần dữ liệu với $x>0$, $y<1/27$; miền từ đoạn dữ liệu còn cần $0\le x/(1-27y)\le1$.

\textbf{Bước 4 — Kiểm tra.} Đặt $D=1-27y$. $u_x=3\sqrt{x}/\sqrt D$, $u_y=27x^{3/2}/D^{3/2}$, nên $u_x^3=u_y$; tại $y=0$ đúng dữ liệu. Ở $x=0$ công thức kéo dài $C^1$ từ phía phải, nhưng $u_{xx}$ không hữu hạn, nên không gọi là nghiệm $C^2$ tại đầu mút này.

%ID BT_1_5_3_11
\textbf{Bước 1 — Gradient đầu.} $F=p+q^2/2-1$, $u(0,s)=s^2$, $0\le s\le1$, cho $q_0=2s$, $p_0=1-2s^2$.

\textbf{Bước 2 — Tích phân.} $p,q$ hằng, $X'=1$, $Y'=q_0=2s$, $U'=p_0+q_0^2=1+2s^2$. Do đó $X=\tau$, $Y=s(1+2\tau)$, $U=s^2+\tau(1+2s^2)$.

\textbf{Bước 3 — Khử.} $\tau=x$, $s=y/(1+2x)$, suy $u=x+y^2/(1+2x)$. Miền cục bộ quanh dữ liệu lấy $1+2x>0$; còn cần $s\in[0,1]$. Công thức đại số cũng có nghĩa trên miền $x<-1/2$, nhưng miền này không nối với dữ liệu qua phép khử đang dùng.

\textbf{Bước 4 — Kiểm tra.} $p=1-2y^2/(1+2x)^2$, $q=2y/(1+2x)$, nên $p+q^2/2=1$; ở $x=0$ bằng $y^2$. Đường $x=-1/2$ làm Jacobian mất khả nghịch.

%ID BT_1_5_3_12
\textbf{Bước 1 — Tương thích.} $F=2xp^2+yq-u$, $u(s,1)=s/2$ cho $p_0=1/2$. Thế PDE được $2s/4+q_0-s/2=0$, nên $q_0=0$.

\textbf{Bước 2 — Giải Charpit.} $p'=p-2p^2$, $q'=0$, nên với dữ liệu này $p=1/2,q=0$. $X'=4Xp=2X$, $Y'=Y$, $U'=4Xp^2+Yq=X$. Suy $X=se^{2\tau}$, $Y=e^\tau$, $U=se^{2\tau}/2$.

\textbf{Bước 3 — Khử và kiểm tra.} $U=X/2$, nên $u=x/2$. Gradient $p=1/2,q=0$ cho $2xp^2+yq=x/2=u$, và vết trên $y=1$ đúng. $J_0=F_q=1$ nên chỉ có một nhánh không đặc trưng cục bộ. Đây là trường hợp hệ có nghiệm đơn giản, không cần giữ mũ ma trận.

%ID BT_1_5_3_13
\textbf{Bước 1 — Kiểm tra tồn tại thực.} $p^2+q^2=u$, dữ liệu $u(s,0)=s^2+1$ cho $p_0=2s$, nên $q_0^2=1-3s^2$. Vì vế phải âm khi $|s|>1/\sqrt3$, không có nghiệm thực trên lân cận toàn trục dữ liệu.

\textbf{Bước 2 — Đổi ẩn trên khoảng hợp lệ.} Gần dữ liệu $u>0$, đặt $v=2\sqrt u$. Vì $v_x=p/\sqrt u$, $v_y=q/\sqrt u$, có $v_x^2+v_y^2=1$ và $v(s,0)=2\sqrt{s^2+1}$. Đặt
\[a(s)=\frac{2s}{\sqrt{s^2+1}},\qquad b(s)=\frac{\pm\sqrt{1-3s^2}}{\sqrt{s^2+1}},\quad a^2+b^2=1.\]

\textbf{Bước 3 — Đặc trưng eikonal.} Dùng tham số độ dài: $X=s+\tau a$, $Y=\tau b$, $v=2\sqrt{s^2+1}+\tau$. Đổi về $u$:
\[u(X,Y)=(\sqrt{s^2+1}+\tau/2)^2.\]
Trên mỗi dấu, khử tham số nơi Jacobian khác 0; tại dữ liệu Jacobian bằng $b$, nên hợp lệ khi $|s|<1/\sqrt3$. Giữ $\sqrt{s^2+1}+\tau/2>0$ cho phép đổi căn.

\textbf{Bước 4 — Kiểm tra.} Công thức cho $v$ có gradient đơn vị trên miền khả nghịch, nên $|\nabla u|^2=(v^2/4)|\nabla v|^2=u$. Ở $\tau=0$ đúng dữ liệu. Hai đầu $|s|=1/\sqrt3$ có $b=0$, cần xét riêng, không áp kết luận không đặc trưng.

\textbf{Chi tiết xác nhận gradient của nghiệm tham số.} Đặt $h=2\sqrt{s^2+1}$, $a=h'$, $a^2+b^2=1$. Với $X=s+\tau a,Y=\tau b,v=h+\tau$, tính $dv=a\,ds+d\tau$. Mặt khác
\[a\,dX+b\,dY=a\,ds+(a^2+b^2)d\tau+\tau(aa'+bb')ds=a\,ds+d\tau,\]
vì đạo hàm $a^2+b^2=1$ cho $aa'+bb'=0$. Trên miền ánh xạ khả nghịch, so hệ số vi phân được $v_x=a,v_y=b$, nên $|\nabla v|^2=1$. Đổi $u=v^2/4$ cho $|\nabla u|^2=(v^2/4)|\nabla v|^2=u$.

%ID BT_1_5_3_14
\textbf{Bước 1 — Đổi ẩn.} Dữ liệu $u(x,0)=1$ dương, nên trong lân cận đủ nhỏ đặt $v=\log u$. Khi đó $v_x=p/u$, $v_y=q/u$ và PDE $p^2+q^2=u^2$ thành $v_x^2+v_y^2=1$, dữ liệu $v(x,0)=0$.

\textbf{Bước 2 — Hai hướng pháp tuyến.} Đạo hàm theo $x$ của vết bằng 0, nên $v_x=0$ trên dữ liệu; PDE cho $v_y=\pm1$. Đặc trưng là đường thẳng pháp tuyến trục $x$, gradient giữ hằng, nên $v=\pm y$.

\textbf{Bước 3 — Đổi lại và thử.} $u=e^y$ hoặc $u=e^{-y}$. Với $p=0,q=\pm u$, có $p^2+q^2=u^2$; trên $y=0$ đều bằng 1. Có hai nhánh cục bộ khi chỉ cho giá trị trên dữ liệu; nếu chọn dấu đạo hàm pháp tuyến thì mỗi nhánh được xác định.

%ID BT_1_5_3_15
\textbf{Bước 1 — Đổi ẩn và dữ liệu.} Đặt $v=\log u$ gần $u=1$. PDE thành $|\nabla v|^2=1$, dữ liệu $v=0$ trên $r=1$.

\textbf{Bước 2 — Chọn gradient ban đầu.} Đạo hàm tiếp tuyến trên vòng tròn bằng 0, nên gradient phải theo pháp tuyến. Độ dài 1 cho hai lựa chọn $\nabla v=\pm(\cos s,\sin s)$. Các đặc trưng là tia pháp tuyến và $v$ thay đổi với tốc độ $\pm1$ theo bán kính, nên $v=\pm(r-1)$.

\textbf{Bước 3 — Nghiệm và kiểm tra.} $u=e^{r-1}$ hoặc $u=e^{1-r}$ trong vành lân cận $r=1$. Vì $|\nabla r|=1$ khi $r>0$, $|\nabla u|^2=(u_r)^2=u^2$; ở $r=1$ bằng 1. Hai nhánh xuyên tâm không kéo dài $C^1$ qua gốc vì $u_r(0)$ khác 0. Dữ liệu không chọn dấu pháp tuyến nên chưa cho duy nhất giữa hai nhánh.

%ID BT_1_5_3_16
\textbf{Bước 1 — Đổi ẩn.} Với $p^2+q^2=u$, gần dữ liệu $u=1>0$, đặt $v=2\sqrt u$; tính gradient như bài 13 cho $|\nabla v|^2=1$, $v|_{r=1}=2$.

\textbf{Bước 2 — Hai nhánh pháp tuyến.} Gradient tiếp tuyến bằng 0, gradient pháp tuyến là $\pm1$, nên $v=2\pm(r-1)$. Đổi lại:
\[u=(r+1)^2/4,\qquad u=(3-r)^2/4.\]

\textbf{Bước 3 — Kiểm tra và miền.} Với $r>0$, $|\nabla u|^2=(u_r)^2$; $u_r=(r+1)/2$ ở nhánh đầu, $(r-3)/2$ ở nhánh sau, nên đều có $(u_r)^2=u$. Tại $r=1$ bằng 1. Phép đổi căn cho nhánh sau dùng $v=3-r>0$, tức $r<3$ quanh dữ liệu; công thức đại số có thể giải PDE xa hơn sau khi kiểm tra trực tiếp. Cả hai nhánh không có kéo dài xuyên tâm $C^1$ qua gốc.

\textbf{Chi tiết đổi ẩn và tia pháp tuyến.} $v=2\sqrt u$ cho $v_x=u_x/\sqrt u,v_y=u_y/\sqrt u$, nên $|\nabla v|^2=(p^2+q^2)/u=1$. Lấy điểm đầu $(\cos s,\sin s)$, gradient $\varepsilon(\cos s,\sin s)$, tham số độ dài $\tau$. Hệ cho $(X,Y)=(1+\varepsilon\tau)(\cos s,\sin s)$, $v=2+\tau$. Gần vòng tròn, $r=1+\varepsilon\tau$, tức $\tau=\varepsilon(r-1)$, nên $v=2+\varepsilon(r-1)$. Bình phương rồi chia $4$ cho hai công thức của $u$.

%ID BT_1_5_3_17
\textbf{Bước 1 — Điều kiện thực trên dữ liệu.} $F=q^2-p$, $u(s,0)=f(s)$ cho $p_0=f'(s)$ và $q_0^2=f'(s)$. Vì vậy cần $f'\ge0$. Nếu $f'<0$ tại một điểm thì vô nghiệm thực quanh điểm đó.

\textbf{Bước 2 — Hai nhánh không đặc trưng.} Khi $f'>0$ và $f$ đủ trơn, đặt $Q=\pm\sqrt{f'}$. $F_p=-1,F_q=2q$, $p'=q'=0$, $U'=-p+2q^2=p$. Tích phân:
\[X=s-\tau,\quad Y=2Q(s)\tau,\quad U=f(s)+f'(s)\tau.\]

\textbf{Bước 3 — Khử và miền.} Giải $x=s-\tau$, $y=2Q(s)\tau$ tìm $s,\tau$, rồi $u=U$. Jacobian là $2Q+2\tau Q'$, tại dữ liệu bằng $2Q\ne0$, nên có nghiệm cục bộ trên từng dấu. Không cần một công thức tường minh theo $x,y$ cho $f$ tùy ý.

\textbf{Bước 4 — Suy biến.} Nếu $f'=0$, dữ liệu đặc trưng và hai dấu trùng nhau. Dữ liệu hằng có nghiệm hằng, nhưng không suy duy nhất từ công thức ở bước 3. Trên miền khả nghịch, Charpit cho gradient $p=f'(s),q=Q(s)$, nên $q^2=p$ và tại $\tau=0$ đúng $f$.

\textbf{Chi tiết phương trình khử và gradient.} Thế $\tau=s-x$ vào $y=2Q(s)\tau$ được một phương trình một ẩn $y=2Q(s)(s-x)$. Đạo hàm theo $s$ là $2Q+2(s-x)Q'$, khác $0$ tại dữ liệu khi $f'>0$, nên giải cục bộ được $s=s(x,y)$. Đặt $P=f'=Q^2$; từ $X=s-\tau,Y=2Q\tau,U=f+P\tau$, tính
\[dU=(P+\tau P')ds+P\,d\tau=P\,dX+Q\,dY,\]
vì $P'=2QQ'$. So hệ số vi phân cho $u_x=P,u_y=Q$, xác nhận $u_y^2=u_x$.

%ID BT_1_5_3_18
\textbf{Bước 1 — Tương thích.} $F=p^2-3q^2-u$, $u(s,0)=s^2$ cho $p_0=2s$. PDE cho $4s^2-3q_0^2=s^2$, nên $q_0=\pm s$.

\textbf{Bước 2 — Tích phân Charpit.} $p'=p,q'=q,X'=2p,Y'=-6q,U'=2U$. Đặt $E=e^\tau$:
\[X=s+4s(E-1),\quad Y=-6q_0(E-1),\quad U=s^2E^2.\]

\textbf{Bước 3 — Khử.} Với $q_0=s$, $X+Y/2=sE$; với $q_0=-s$, $X-Y/2=sE$. Vì thế $u=(x+y/2)^2$ hoặc $u=(x-y/2)^2$.

\textbf{Bước 4 — Kiểm tra.} Đặt $L=x\pm y/2$. $p=2L,q=\pm L$, nên $p^2-3q^2=L^2=u$, và $y=0$ cho $x^2$. $J_0=-6q_0=\mp6s$ khác 0 khi $s\ne0$; gốc là điểm suy biến.

%ID VD_1_2_2_1_01
\textbf{Bước 1 — Xác định cấp.} Liệt kê $u_{xxx}$ cấp 3, các đạo hàm còn lại có cấp thấp hơn. Bình phương $(u_y)^2$ vẫn chỉ chứa đạo hàm cấp 1. Do đó cấp cao nhất là $m=3$.

\textbf{Bước 2 — Xét phần chính.} Đạo hàm $u_{xxx}$ xuất hiện ở lũy thừa 1 với hệ số 1, không phụ thuộc $u$ hay đạo hàm của $u$.

\textbf{Bước 3 — Xét phần cấp thấp.} Hạng $(u_y)^2$ phi tuyến nên toàn phương trình không tuyến tính. Vì phi tuyến chỉ ở cấp thấp và hệ số phần chính chỉ phụ thuộc biến độc lập, kết luận: \textbf{nửa tuyến tính, cấp 3}.

%ID VD_1_2_2_1_02
\textbf{Bước 1 — Đếm cấp.} Đạo hàm cao nhất xuất hiện là $u_x$, nên $m=1$; $u^3$ là hạng không có đạo hàm.

\textbf{Bước 2 — Xét tính tuyến tính theo đạo hàm cấp 1.} Coi $p=u_x$ như một biến đại số. Biểu thức lượng giác và bình phương chứa $p$ không có dạng $A(x,y,u)p+B(x,y,u)$. Như vậy phi tuyến xuất hiện ngay ở đạo hàm cấp cao nhất.

\textbf{Kết luận.} Phương trình \textbf{hoàn toàn phi tuyến, cấp 1}. Hạng $u^3$ không làm tăng cấp và không phải lý do chính để phân biệt hoàn toàn phi tuyến với nửa tuyến tính.

%ID VD_1_2_2_1_03
\textbf{Bước 1 — Xác định cấp.} Các đạo hàm $u_{xx},u_{xy},u_{yy}$ đều có cấp 2, nên $m=2$.

\textbf{Bước 2 — Đọc các hệ số.} Các hệ số tương ứng là $x,x-y,-y$, chỉ phụ thuộc $x,y$. Mỗi đạo hàm xuất hiện tuyến tính; không có tích giữa các đạo hàm hoặc tích với ẩn $u$.

\textbf{Bước 3 — Xét nguồn.} Phương trình đã có dạng $Lu=0$, nên thuần nhất. Kết luận: \textbf{tuyến tính thuần nhất, cấp 2}. Hệ số thay đổi theo vị trí không làm mất tính tuyến tính.

%ID VD_1_2_2_1_04
\textbf{Bước 1 — Xác định cấp.} $u_{xx}$ là đạo hàm cao nhất, nên $m=2$.

\textbf{Bước 2 — Xét dạng của phần cấp 2.} Đạo hàm $u_{xx}$ xuất hiện tuyến tính, nhưng hệ số là $u^2$, phụ thuộc chính ẩn cần tìm. Vì thế không thể xem phương trình là tuyến tính hoặc nửa tuyến tính theo định nghĩa trong sách.

\textbf{Kết luận.} Phương trình \textbf{tựa tuyến tính, cấp 2}. Tại $u=0$, hệ số phần chính có thể triệt tiêu; cần lưu ý sự suy biến khi giải, nhưng phân loại biểu thức chung vẫn như trên.

%ID VD_1_2_2_2_01
\textbf{Bước 1 — Chọn biến tích phân.} Phương trình là $u_x=3x^2+y^2$. Giữ $y$ cố định và tích phân theo $x$.

\textbf{Bước 2 — Tính từng tích phân.}
\[u=\int3x^2\,dx+\int y^2\,dx+C(y)=x^3+xy^2+C(y).\]
Hàm $C(y)$ xuất hiện vì đạo hàm theo $x$ của mọi hàm chỉ phụ thuộc $y$ đều bằng 0.

\textbf{Bước 3 — Kiểm tra.} $u_x=3x^2+y^2$ đúng phương trình. Không có dữ liệu bổ sung để xác định $C$; chọn $C$ đủ trơn theo lớp nghiệm yêu cầu.

%ID VD_1_2_2_2_02
\textbf{Bước 1 — Giảm đạo hàm hỗn hợp.} Từ $u_{xyz}=0$, cố định $y,z$ và tích phân theo $x$, được $u_{yz}=a(y,z)$.

\textbf{Bước 2 — Tích phân theo $y$, rồi $z$.} Chọn một nguyên hàm $b(y,z)$ của $a$ theo $y$, được $u_z=b(y,z)+c(x,z)$. Tích phân theo $z$:
\[u(x,y,z)=B(y,z)+C(x,z)+A(x,y).\]
Mỗi hàm tùy ý bỏ một biến; điều này giữ đủ các hằng số tích phân của ba lần lấy nguyên hàm. Công thức có ý nghĩa trên miền tích hoặc cục bộ nơi các phép tích phân được thực hiện.

\textbf{Bước 3 — Kiểm tra.} $A$ không phụ thuộc $z$, $B$ không phụ thuộc $x$, $C$ không phụ thuộc $y$, nên đạo hàm $xyz$ của từng hạng bằng 0. Vì vậy tổng thỏa PDE.

%ID VD_1_3_1_2_03
\textbf{Bước 1 — Lập đặc trưng.} $u_t-2u_x=2$ cho $X'=-2$, $U'=2$ với $t$ là tham số. Từ $X+2t=\xi$, có bất biến $x+2t$.

\textbf{Bước 2 — Tích phân ẩn.} $U=2t+F(\xi)$, nên $u=2t+F(x+2t)$. Cũng có thể chọn nghiệm riêng $2t$, rồi giải PDE thuần nhất cho phần còn lại.

\textbf{Bước 3 — Kiểm tra.} $u_t=2+2F'$, $u_x=F'$, nên $u_t-2u_x=2$. Không có dữ liệu để xác định $F$.

%ID VD_1_3_1_2_04
\textbf{Bước 1 — Đường đặc trưng.} Với $u_x+2u_y+u=0$, dùng $x$ làm tham số: $dy/dx=2$, nên $\rho=y-2x$ không đổi.

\textbf{Bước 2 — Tích phân ẩn.} Dọc đường, $dU/dx+U=0$, nên $(e^xU)'=0$. Vì vậy $u=e^{-x}F(y-2x)$.

\textbf{Bước 3 — Dùng $u(x,0)=x^2$.} $F(-2x)=x^2e^x$. Đặt $s=-2x$, được $F(s)=s^2e^{-s/2}/4$; suy ra
\[u(x,y)=\frac{(y-2x)^2}{4}e^{-y/2}.\]

\textbf{Bước 4 — Kiểm tra.} Trên $y=0$ công thức bằng $x^2$. Với dạng tổng quát, $u_x=e^{-x}(-F-2F')$, $u_y=e^{-x}F'$, nên $u_x+2u_y+u=0$. Mỗi đặc trưng gặp $y=0$ đúng một lần, nên nghiệm được xác định duy nhất.

%ID VD_1_3_1_2_05
\textbf{Bước 1 — Tìm bất biến.} Với $2u_x-3u_y=x$, $dy/dx=-3/2$, nên $\rho=3x+2y$.

\textbf{Bước 2 — Tích phân trên đường.} $dU/dx=x/2$, do đó $U=x^2/4+F(\rho)$ và
\[u=x^2/4+F(3x+2y).\]

\textbf{Bước 3 — Kiểm tra.} $u_x=x/2+3F'$, $u_y=2F'$, nên $2u_x-3u_y=x+6F'-6F'=x$.

%ID VD_1_3_1_2_06
\textbf{Bước 1 — Đường đặc trưng.} $u_x+2u_y-4u=e^{x+y}$ có $dy/dx=2$, nên $\rho=y-2x$. Trên một đường $y=2x+\rho$, ẩn thỏa $U'-4U=e^{3x+\rho}$.

\textbf{Bước 2 — Nhân thừa số tích phân.}
\[(e^{-4x}U)'=e^{\rho-x},\quad e^{-4x}U=-e^{\rho-x}+F(\rho).\]
Suy ra $u=-e^{x+y}+e^{4x}F(y-2x)$. Dạng này tương đương dạng mũ $e^{4(x+y)/3}$ trong sách sau khi đổi tên hàm tùy ý.

\textbf{Bước 3 — Xác định $F$ từ dữ liệu.} $u(x,0)=\sin(x^2)$ cho
\[F(-2x)=e^{-4x}[\sin(x^2)+e^x].\]
Đặt $s=-2x$, được $F(s)=e^{2s}\sin(s^2/4)+e^{3s/2}$, nên
\[u=e^{2y}\sin((y-2x)^2/4)+e^{x+3y/2}-e^{x+y}.\]

\textbf{Bước 4 — Kiểm tra.} Với $L=\partial_x+2\partial_y-4$, $L(e^{4x}F(\rho))=0$, còn $L(-e^{x+y})=e^{x+y}$. Trên $y=0$, hai hạng mũ triệt tiêu và còn $\sin(x^2)$. Hệ số theo $y$ khác 0 nên mỗi đặc trưng gặp dữ liệu đúng một lần.

%ID VD_1_3_1_2_07
\textbf{Bước 1 — Giải PDE.} $u_x+u_y-u=0$ cho $dy/dx=1$, $\rho=y-x$ và $U'=U$, nên $u=e^xF(y-x)$.

\textbf{Bước 2 — Thế dữ liệu trên đặc trưng.} Trên $y=x$, $\rho=0$, nên $u(x,x)=e^xF(0)$. Để bằng $\tan x$, phải có $F(0)=e^{-x}\tan x$ ở mọi điểm của khoảng dữ liệu.

\textbf{Bước 3 — Chứng minh mâu thuẫn.} Ở $x=0$, điều kiện cho $F(0)=0$, nhưng ở $x$ gần 0 khác 0, $e^{-x}\tan x\ne0$. Do đó không có nghiệm trên lân cận đường dữ liệu quanh gốc. Cũng có thể dùng điều kiện cần $g'=g$ trên đường; $g=\tan x$ không thỏa.

%ID VD_1_3_1_2_08
\textbf{Bước 1 — Dùng nghiệm tổng quát.} Như ví dụ trước, $u=e^xF(y-x)$. Trên $y=x$, chỉ số đặc trưng bằng 0, nên dữ liệu phải là $f(x)=e^xF(0)$.

\textbf{Bước 2 — Điều kiện cần và đủ.} Đặt $A=F(0)$, ta có $f(x)=Ae^x$. Ngược lại, nếu $f$ có dạng này, mọi hàm $F$ đủ trơn thỏa $F(0)=A$ đều cho nghiệm.

\textbf{Bước 3 — Giải thích không duy nhất.} Chọn $F_1(s)=A$ và $F_2(s)=A+s$ cho $u_1=Ae^x$, $u_2=e^x(A+y-x)$. Hai hàm cùng thỏa PDE và dữ liệu, nhưng khác nhau ngoài $y=x$. Dữ liệu trên một đặc trưng chỉ xác định giá trị của $F$ tại một điểm.

\textbf{Chi tiết dựng nghiệm tổng quát trước khi dùng vết.} Với $u_x+u_y=u$, chọn $x$ làm tham số: $Y'=1$, $U'=U$. Tích phân $Y=x+\rho$ và $(e^{-x}U)'=0$ cho $e^{-x}U=F(\rho)$. Vì $\rho=y-x$, thu $u=e^xF(y-x)$. Như vậy công thức dùng ở bước 1 xuất phát từ chính hệ đặc trưng của bài này.

%ID VD_1_3_1_2_09
\textbf{Bước 1 — Giải theo $x$.} Với $au_x+bu_y+cu=0$, $ab\ne0$, bất biến $\rho=bx-ay$ và $dU/dx+(c/a)U=0$. Suy ra $u=e^{-cx/a}F(\rho)$.

\textbf{Bước 2 — Giải theo $y$.} Tương tự $dU/dy+(c/b)U=0$, nên $u=e^{-cy/b}G(\rho)$. Cả hai cách đều đúng, với hàm tùy ý tương ứng.

\textbf{Bước 3 — Chuyển giữa hai biểu thức.} Vì $x/a-y/b=(bx-ay)/(ab)=\rho/(ab)$,
\[e^{-cx/a}F(\rho)=e^{-cy/b}e^{-c\rho/(ab)}F(\rho).\]
Do đó đặt $G(s)=e^{-cs/(ab)}F(s)$ để được cùng nghiệm. Không cần $F=G$; hai bạn đang dùng hai cách tham số hóa cùng họ nghiệm.

%ID VD_1_3_2_2_02
\textbf{Ý (a) — $u_x+x^2u_y=0$.} Bước 1: $dy/dx=x^2$ cho $y=x^3/3+\rho$. Bước 2: $U'=0$, nên $u=F(y-x^3/3)$. Bước 3: $u_x=-x^2F'$, $u_y=F'$, thế vào PDE được 0.

\textbf{Ý (b) — $e^{x^2}u_x+xu_y=0$.} Vì $e^{x^2}>0$, chọn $x$ làm tham số. $dy/dx=xe^{-x^2}$; dùng $d(e^{-x^2})=-2xe^{-x^2}dx$, suy $y=-e^{-x^2}/2+\rho/2$. Do đó $u=F(2y+e^{-x^2})$. Kiểm tra $u_x=-2xe^{-x^2}F'$, $u_y=2F'$, nên $e^{x^2}u_x+xu_y=0$. Cả hai công thức hợp lệ trên toàn mặt phẳng với $F\in C^1$.

%ID VD_1_3_2_2_04
\textbf{Bước 1 — Tìm đặc trưng.} $xu_x-2yu_y+u=e^x$, $x>0$. Hệ $X'=X,Y'=-2Y$ cho $(X^2Y)'=0$, nên $\rho=x^2y$. Cách tính này bao gồm $y=0$, không cần lấy log của $y$.

\textbf{Bước 2 — Giải ODE.} Dùng $x$ làm tham số, $U'+U/x=e^x/x$. Vì $x>0$, thừa số tích phân là $x$:
\[(xU)'=e^x,\quad xU=e^x+F(\rho).\]
Suy $u=[e^x+F(x^2y)]/x$.

\textbf{Bước 3 — Kiểm tra.} Đặt $V=x\partial_x-2y\partial_y$. $V\rho=0$, $V(1/x)=-1/x$; do đó $(V+1)(F(\rho)/x)=0$. Với $e^x/x$, $(V+1)(e^x/x)=e^x$. Công thức thỏa PDE trên miền $x>0$.

%ID VD_1_3_2_2_05
\textbf{Bước 1 — Dùng nghiệm tổng quát của ví dụ 4.} $u=[e^x+F(x^2y)]/x$ trên $x>0$.

\textbf{Bước 2 — Thế dữ liệu $u(1,y)=y$.} $e+F(y)=y$, nên $F(s)=s-e$. Thay lại:
\[u(x,y)=xy+(e^x-e)/x.\]

\textbf{Bước 3 — Kiểm tra và duy nhất.} Trên $x=1$, công thức bằng $y$. Phần $xy$ thỏa phương trình thuần nhất vì $x(xy)_x-2y(xy)_y+xy=0$; phần còn lại cho nguồn $e^x$. Mỗi đặc trưng $x^2y=\rho$ trong $x>0$ gặp $x=1$ đúng một lần, nên duy nhất trên miền này.

\textbf{Chi tiết giải lại PDE.} Hệ $X'=X,Y'=-2Y$ cho $X=x_0e^s,Y=y_0e^{-2s}$, nên $X^2Y=x_0^2y_0=\rho$. Vì $dX/ds=X>0$, phương trình của ẩn theo $x$ là $dU/dx+U/x=e^x/x$. Nhân $x$ rồi tích phân: $d(xU)/dx=e^x$, $xU=e^x+F(\rho)$. Do đó $u=(e^x+F(x^2y))/x$, trước khi thế $x=1$ như trong bước 2.

%ID VD_1_3_2_2_06
\textbf{Bước 1 — Nhận ra đường dữ liệu là đặc trưng.} Trong nghiệm $u=[e^x+F(x^2y)]/x$, trên $y=0$ chỉ số $x^2y=0$ không đổi.

\textbf{Bước 2 — Tìm điều kiện tương thích.} $u(x,0)=f(x)$ đòi $xf(x)-e^x=F(0)=C_0$, tức $f(x)=(e^x+C_0)/x$ trên $x>0$. Đây là điều kiện cần và đủ.

\textbf{Bước 3 — Chỉ ra vô số nghiệm.} Nếu dữ liệu tương thích, chọn $F_\lambda(s)=C_0+\lambda s$ với mọi $\lambda\in\mathbb R$. Khi đó
\[u_\lambda=(e^x+C_0)/x+\lambda xy.\]
Các nghiệm đều thỏa PDE và vết $f$, nhưng khác nhau ngoài $y=0$. Có thể dùng mọi $F\in C^1$ với $F(0)=C_0$; họ vừa viết đủ chứng minh không duy nhất.

\textbf{Chi tiết nguồn gốc của hàm tùy ý.} Dọc $X'=X,Y'=-2Y$, tính $(X^2Y)'=2X^2Y-2X^2Y=0$. Phương trình $dU/dx+U/x=e^x/x$ cho $(xU)'=e^x$, rồi $xU=e^x+F(X^2Y)$. Đường $y=0$ chỉ gặp chỉ số $\rho=0$, nên nó chỉ quyết định $F(0)$; các giá trị $F(\rho)$ với $\rho\ne0$ không được quyết định.

%ID VD_1_4_2_03
\textbf{Bước 1 — Thông lượng và đặc trưng.} Burgers có $F=u^2/2$. Dữ liệu $u_L=1,u_R=2$ ở $x=0$ cho $X=\xi+t$ từ trái, $X=\xi+2t$ từ phải. Hai họ tách xa, để trống $t<x<2t$.

\textbf{Bước 2 — Lấp bằng quạt giãn.} Đặt $z=x/t$, $u=V(z)$. PDE cho $(V-z)V'=0$, nên nhánh biến thiên có $V=z$. Nó ghép liên tục với giá trị 1 tại $x=t$ và 2 tại $x=2t$.

\textbf{Bước 3 — Nghiệm và kiểm tra.}
\[u=\begin{cases}1,&x\le t,\\x/t,&t<x<2t,\\2,&x\ge2t,\end{cases}\quad t>0.\]
Trong quạt $u_t=-x/t^2$, $uu_x=x/t^2$, tổng bằng 0. Các mép liên tục và dữ liệu khi $t\downarrow0$ đúng ở $x\ne0$. Đây là nghiệm entropy; dữ liệu gián đoạn không cho nghiệm cổ điển qua $t=0$.

%ID VD_1_5_2_01
\textbf{Ý (a), bước 1 — Tương thích dữ liệu.} Đặt $F=p^2+q^2-1$, $X(s,0)=s,Y(s,0)=0,U(s,0)=f(s)$. Đạo hàm vết cho $p_0=f'(s)$, còn PDE cho $q_0^2=1-f'(s)^2$. Nếu $|f'(0)|>1$, không có nghiệm thực gần gốc.

\textbf{Bước 2 — Không đặc trưng và hệ Charpit.} Nếu $|f'(0)|<1$ và $f\in C^2$ gần gốc, thu hẹp khoảng để $|f'(s)|<1$. Có hai nhánh $q_0=\pm\sqrt{1-f'(s)^2}$ và $J_0=F_q=2q_0\ne0$. Hệ là
\[X'=2p,\quad Y'=2q,\quad U'=2,\quad p'=q'=0.\]
Do đó
\[X=s+2\tau f'(s),\quad Y=2\tau q_0(s),\quad U=f(s)+2\tau.\]
Trên từng nhánh, khử $s,\tau$ gần dữ liệu để được một nghiệm cục bộ. Nếu $f(s)=s/2$, khử trực tiếp cho $u=x/2\pm\sqrt3\,y/2$.

\textbf{Bước 3 — Trường hợp $|f'(0)|=1$.} $J_0=0$, nên không dùng định lý không đặc trưng để kết luận duy nhất. Dữ liệu $f(x)=x+x^2\operatorname{sgn}x$ có $f'(x)=1+2|x|>1$ khi $x\ne0$ gần gốc, nên vô nghiệm dù $f'(0)=1$. Dữ liệu $f(x)=x$ lại có vô số nghiệm $u_C=C-\sqrt{(x-C)^2+y^2}$ gần gốc, với $C>1$: trên $y=0,x<C$ có $u_C=x$, và $|\nabla u_C|=1$. Các ví dụ cho thấy phải xét cả dữ liệu lân cận, không chỉ giá trị đạo hàm tại gốc.

\textbf{Ý (b), bước 1 — Dữ liệu trên vòng tròn.} $x_0=\cos s,y_0=\sin s,h=0$. Đạo hàm tiếp tuyến cho $-p_0\sin s+q_0\cos s=0$; PDE cho $p_0^2+q_0^2=1$. Vì vậy gradient có hai lựa chọn $(p_0,q_0)=\varepsilon(\cos s,\sin s)$, $\varepsilon=\pm1$.

\textbf{Bước 2 — Tích phân và khử.} $X=(1+2\varepsilon\tau)\cos s$, $Y=(1+2\varepsilon\tau)\sin s$, $U=2\tau$. Gần vòng tròn giữ $1+2\varepsilon\tau>0$, nên $r=1+2\varepsilon\tau$ và
\[u=\varepsilon(r-1),\qquad r=\sqrt{x^2+y^2}.\]

\textbf{Bước 3 — Kiểm tra và sửa kết luận trong sách.} Hai hàm $r-1$ và $1-r$ đều có $|\nabla u|=1$ và vết 0 tại $r=1$. $J_0=-2\varepsilon\ne0$ cho duy nhất cục bộ \emph{trong từng nhánh gradient}, nhưng dữ liệu giá trị chưa chọn dấu nên có hai nhánh. Không kéo dài hai hàm xuyên tâm thành nghiệm $C^1$ qua gốc.

\textbf{Chi tiết khử tham số và chọn gradient.} Ở (a) với $f=s/2$, $q_0=\varepsilon\sqrt3/2$, $X=s+\tau$, $Y=\varepsilon\sqrt3\tau$, $U=s/2+2\tau$. Do đó $\tau=\varepsilon y/\sqrt3$, $s=x-\tau$, $U=x/2+3\tau/2=x/2+\varepsilon\sqrt3 y/2$. Ở (b), đặt tiếp tuyến $e_\theta=(-\sin s,\cos s)$ và pháp tuyến $e_r=(\cos s,\sin s)$. Viết $(p_0,q_0)=Ae_r+Be_\theta$; dữ liệu cho $B=0$, PDE cho $A^2+B^2=1$, nên $A=\varepsilon$. Từ $r=1+2\varepsilon\tau$ có $\tau=\varepsilon(r-1)/2$, $U=2\tau=\varepsilon(r-1)$.

%ID VD_1_5_2_02
\textbf{Bước 1 — Tương thích.} $F=p^2+q^2+2(p-x)(q-y)-2u$, dữ liệu $(s,0,0)$ cho $p_0=0$ và $q_0(q_0-2s)=0$. Có hai nhánh $q_0=0,2s$; $J_0=2(q_0-s)$ khác 0 khi $s>0$.

\textbf{Bước 2 — Hệ và các bất biến.} Tính đạo hàm của $F$:
\[X'=2(p+q-Y),\quad Y'=2(p+q-X),\quad p'=X',\quad q'=Y'.\]
Do đó $X-p=s$, $Y-q=-q_0$. Thế $p=X-s,q=Y+q_0$ vào hệ đường đi, được $X'=2X+2(q_0-s)$, $Y'=2Y+2(q_0-s)$. Tích phân với dữ liệu:
\[X=q_0e^{2\tau}+s-q_0,\qquad Y=(q_0-s)(e^{2\tau}-1).\]
Từ $F=0$,
\[U=\frac{X^2+Y^2+s^2-2Xs+2Yq_0-2sq_0+q_0^2}{2}.\]

\textbf{Bước 3 — Khử hai nhánh.} Nếu $q_0=0$, $X=s$, $Y=s(1-e^{2\tau})$, suy $p=0,q=Y$ và $u=y^2/2$. Nếu $q_0=2s$, $X=2se^{2\tau}-s$, $Y=s(e^{2\tau}-1)$, nên $s=X-2Y$, $p=2Y,q=2X-3Y$, suy $u=2xy-3y^2/2$.

\textbf{Bước 4 — Kiểm tra và miền.} Với $u=y^2/2$, $p=0,q=y$ cho vế trái $y^2=2u$. Với $u=2xy-3y^2/2$, $p=2y,q=2x-3y$ cho vế trái $4xy-3y^2=2u$. Cả hai có $u(x,0)=0$. Từ đoạn dữ liệu $0\le s\le1$, nhánh thứ nhất cần $s=x$ và $e^{2\tau}=1-y/x>0$; nhánh thứ hai cần $s=x-2y$ và $e^{2\tau}=(x-y)/(x-2y)>0$. Các đa thức là nghiệm kéo dài toàn cục, nhưng dữ liệu chỉ bảo đảm duy nhất cục bộ trên mỗi nhánh ở $s>0$; gốc suy biến. Lời giải sách nói duy nhất khi đã chọn gradient ban đầu; nếu chưa chọn nhánh thì có hai nghiệm.

\textbf{Chi tiết giải hai ODE tuyến tính.} Đặt $c=q_0-s$, hằng theo $\tau$. Từ $X'-2X=2c$, nhân $e^{-2\tau}$ và tích phân từ $0$:
\[e^{-2\tau}X-s=c(1-e^{-2\tau}),\quad X=(s+c)e^{2\tau}-c.\]
Tương tự $e^{-2\tau}Y=c(1-e^{-2\tau})$, nên $Y=c(e^{2\tau}-1)$. Vì $s+c=q_0$, đây đúng các công thức ở bước 2. Nhánh $q_0=2s$ cho $X-2Y=s$, $p=X-s=2Y$, $q=Y+2s=2X-3Y$. Tích phân $u_x=2y$ cho $u=2xy+C(y)$; phương trình $u_y=2x-3y$ cho $C'=-3y$, rồi $C=-3y^2/2$ nhờ vết ban đầu.

%ID VD_1_6_2_01
\textbf{Bước 1 — Đặc trưng và vùng giao.} $F=u^2/2$, dữ liệu 1 khi $\xi<0$, 0 khi $\xi>0$. Vì $U'=0$, $X=\xi+t$ từ trái và $X=\xi$ từ phải. Vùng $x<0$ chỉ có đặc trưng trái, $x>t$ chỉ có đặc trưng phải; vùng $0<x<t$ hai họ giao nhau.

\textbf{Bước 2 — Tìm sốc.} Nối $u_L=1,u_R=0$, Rankine–Hugoniot cho $s'=(1/2-0)/(1-0)=1/2$. Với $s(0)=0$, $s=t/2$.

\textbf{Bước 3 — Nghiệm entropy.} $u=1$ khi $x<t/2$, $u=0$ khi $x>t/2$. Sốc nằm trong vùng giao $0<x<t$ và thỏa $F'(0)=0<s'<F'(1)=1$.

\textbf{Bước 4 — Đọc đồ thị.} Trong mặt phẳng $(x,t)$, vẽ họ $x=\xi+t$ ở trái, họ đứng $x=\xi$ ở phải và đường $x=t/2$ phân chia; đặc trưng đi vào sốc từ cả hai phía. Đồ thị theo $x$ ở mỗi thời điểm là bước nhảy 1 xuống 0 tại $t/2$. Giá trị trên chính sốc tùy chọn không đổi nghiệm yếu.

%ID VD_1_6_2_02
\textbf{Bước 1 — Đặc trưng tách xa.} Dữ liệu Burgers 0 khi $\xi<0$, 1 khi $\xi>0$ cho $X=\xi$ từ trái, $X=\xi+t$ từ phải. Vùng $0<x<t$ không có đặc trưng từ một điểm dữ liệu trơn.

\textbf{Bước 2 — Bổ sung quạt.} Tìm $u=V(z)$, $z=x/t$. Tính $u_t=-zV'/t$, $u_x=V'/t$, nên $(V-z)V'=0$. Nhánh nối liên tục 0 với 1 có $V=z$.

\textbf{Bước 3 — Công thức và kiểm tra.}
\[u=\begin{cases}0,&x\le0,\\x/t,&0<x<t,\\1,&x\ge t,\end{cases}\quad t>0.\]
PDE đúng trong quạt và hai vùng hằng; giá trị ghép tại mép là 0,1, nên không có bước nhảy cần tính sốc.

\textbf{Bước 4 — Đọc đồ thị.} Vẽ các tia $x=kt$, $0\le k\le1$ từ gốc để lấp vùng giữa hai họ ngoài. Trên mỗi tia $u=k$. Ở thời điểm cố định, đồ thị có đoạn thẳng từ $(0,0)$ đến $(t,1)$ giữa hai đoạn hằng. Giới hạn tại $t=0$ khớp dữ liệu ở $x\ne0$.

%ID VD_1_6_2_03
\textbf{Bước 1 — Sóng ban đầu.} Burgers có xung 1 trên $(0,1)$, 0 ở ngoài. Bước $0\to1$ tại 0 tạo quạt $u=x/t$ trên $0<x<t$; bước $1\to0$ tại 1 tạo sốc $s_1=1+t/2$.

\textbf{Bước 2 — Thời điểm sốc gặp quạt.} Giải $1+t/2=t$, được $t_*=2,x_*=2$. Trước đó,
\[u=\begin{cases}0,&x\le0,\\x/t,&0<x<t,\\1,&t\le x<1+t/2,\\0,&x>1+t/2.\end{cases}\quad0<t<2.\]

\textbf{Bước 3 — Sốc trong quạt.} Sau $t=2$, phía trái sốc là $u_L=s/t$, phải là 0. Rankine–Hugoniot cho $s'=s/(2t)$. Tích phân $d\log s=dt/(2t)$ cho $s=C\sqrt t$; dùng $s(2)=2$, suy $s=\sqrt{2t}$. Với $t\ge2$, $u=x/t$ trên $0<x<s$, bằng 0 ở ngoài.

\textbf{Bước 4 — Kiểm tra và hình học.} Sốc ghép liên tục ở 2, tốc độ sau gặp bằng nửa trạng thái trái nên thỏa entropy. Khối lượng sau gặp $\int_0^s(x/t)dx=s^2/(2t)=1$, đúng xung ban đầu. Đồ thị có quạt thẳng từ gốc và bước nhảy xuống 0 trên đường cong $x=\sqrt{2t}$; không kéo dài đường thẳng $1+t/2$ sau tương tác.
